%% 第31回編集委員会（物理第5巻）：組 原W1
%% p5-02.tex  原-2　光はつぶでもある：光電効果と光子（★2〜5・3段：N-1）
%% 中身は計画書 tools/phys5-plan.md の B「原-2」。公式の持ち主は A（A17・A18・A44）。コンプトン効果（A22）は⑦で式を読むだけ（導くのは原-11 ⑧）。
%% 3段：目でわかる（①②、★2）→ 手でわかる（③〜⑧、★3。ここまでで十分）→ 式でわかる（⑨〜⑪ ★5、⑫は工学）。
%%   X-1 の物差しで、プランクの式をモードの数え上げで導くことは★5（計画書 B のはじめ）。手でわかるの段は、高校の E＝hf・K₀＝hf−W・p＝h/λ だけでたどれる。
%% 前の巻の約束（計画書 A-6）：Q5・Q22（プランク定数 h：③、記号カード）、Q9（熱-13 の c₁・c₂ を h で：⑪・例題12）、Q16（ボース分布：⑩）、
%%   Q23（太陽の表面の温度：①）、Q26（1900年のプランク：③）、Q32・Q39（光は粒：③）、Q36（光子を1つずつ数える干渉じま：⑧）、Q59（光子の運動量
%%   ＝エネルギー÷c、1361 W/m² で 4.5 μPa：⑥・例題7）、Q61・Q62（電磁波がつぶとしても：③⑥）。熱-15 ⑧の ⟨ε⟩＝Δε/(e^{Δε/kT}−1) を⑩で Δε＝hf として使った（Q11 の
%%   「比例の係数は原-2・原-6」の原-2 の分）。熱-13 の c₁＝3.742×10⁸ W·μm⁴/m²・c₂＝14388 μm·K・5772 K・可視光 約47 % とそろえた。
%% 数学との約束（C-6-2）：M7（7-10 の S14 等比級数を x＝e^{−a} の形で：⑩）。5-10（台形の和：例題1）、5-7（山の位置：⑪）、5-14（e^x）を指した。
%% 数値はすべて Python で計算した：checks/物理第5巻/原-2/calc_p5_02.py（出力 calc_p5_02_out.txt）。例題1の囲みのプログラムは sim_visible.py（紙面と同じ）。
%% 記号（この章）：プランク定数 h、光の振動数 f（高校の ν：N-5）、波長 λ、光子のエネルギー E、運動量 p、仕事関数 W（この章では仕事を記号にしない）、
%%   飛び出す電子の運動エネルギーの最大値 K₀、阻止電圧 V₀、限界振動数 f₀、電気素量 e、電子の質量 m（m_e）、光速 c、温度 T、ボルツマン定数 k_B、
%%   黒体の単色放射の強さ E_bλ（熱-13 と同じ）、箱の1辺 L、モードの数 N、モードを数える整数 n_x・n_y・n_z、x＝hf/(k_B T)（または hc/(λk_B T)）、
%%   1 m³・1 Hz あたりの光のエネルギー u_f、モード1つの平均のエネルギー ⟨ε⟩（熱-15 と同じ）、シュテファン・ボルツマン定数 σ、ウィーンの定数 b_W（ねらいの距離 b：原-1 とぶつけない）、
%%   散乱角 θ、電子のコンプトン波長 λ_C。
%% 出典（11月11日に確かめた）
%%   ・定数：CODATA 2022（NIST allascii.txt を curl で読んだ）：h（定義値）・e・c・k_B（定義値）、mₑ、電子のコンプトン波長 2.42631023538 pm、
%%       σ＝5.670374419×10⁻⁸、ウィーンの定数 2.897771955×10⁻³ m·K、c₁＝3.741771852×10⁻¹⁶ W·m²、c₂＝1.438776877×10⁻² m·K。
%%   ・ノーベル賞（nobelprize.org、計画書 D15・D16・D18・D19・D38）：プランク 1918年（エネルギー量子）、アインシュタイン 1921年（光電効果の法則）、
%%       ミリカン 1923年（電気素量と光電効果）、コンプトン 1927年、アゴスティーニ・クラウス・ルイリエ 2023年（アト秒の光のパルス）。
%%   ・ミリカンの光電効果の測定：R. A. Millikan, "A direct photoelectric determination of Planck's h", Phys. Rev. 7, 355 (1916)（APS・AIP の資料で照らし合わせた）。
%%   ・光電子が飛び出す時間の差：M. Schultze ほか "Delay in Photoemission", Science 328, 1658 (2010)、ネオンの 2p と 2s の電子で 21±5 アト秒（MPQ の発表 2010年6月25日、
%%       TUM の書誌で照らし合わせた）。差の理由は、原子の中の電子どうしの影響だという計算（Kheifets と Ivanov, 2010）があるが、議論が続いている。
%%   ・モリブデンの Kα1 線 0.709 317 Å（70.93 pm）：NIST X-ray Transition Energies Database（physics.nist.gov/cgi-bin/XrayTrans を curl で読んだ）。
%%   ・セシウム137（バリウム137m）のγ線 661.657 keV：二次資料2つ（ORNL の崩壊図 epa-sdcc.ornl.gov/Cs-137DecayFig.pdf の抜粋、Ionactive の安全資料）で照らし合わせた（2026年11月）。
%%       DDEP・NNDC は作業環境から届かない（計画書 D の記録と同じ）。本文は 0.662 MeV。
%%   ・仕事関数（目安）：Cs 約2.1、K 約2.3、W 約4.5、Cu 約4.7 eV。HyperPhysics の表（Georgia State Univ.）と、H. B. Michaelson, J. Appl. Phys. 48, 4729 (1977) の値
%%       （検索の抜粋）で照らし合わせた。表面の状態で 0.1〜0.5 eV ちがうので、本文は「目安」と書いた。理科年表は引いていない。
%%   ・禁制帯の幅（300 K）：ケイ素 1.12 eV、ゲルマニウム 0.66 eV。S. M. Sze "Physics of Semiconductor Devices" の付録の値を引く資料（BYU Cleanroom の表、
%%       Ioffe 研究所の NSM のミラー www-f9.ijs.si）で照らし合わせた。
%%   ・太陽：5772 K・1361 W/m²（IAU 2015 Resolution B3、計画書 D46）。熱-13・電-15 と同じ。
%%   ・テイラーの弱い光の干渉じま：G. I. Taylor, "Interference fringes with feeble light", Proc. Camb. Phil. Soc. 15, 114 (1909)。
%% 〔要確認〕なし。例題の仮の値（例題5の阻止電圧のデータ、LED の 1 mW など）は本文に「〜とします」と書いた。
\makeatletter
\tikzset{txc lab/.style={font=\sffamily\footnotesize, text=bkInk},
  txc ph/.style={bkOrange!85!black, line width=1.1pt, decorate, decoration={snake, amplitude=1.1mm, segment length=3.2mm, post length=2.4mm}, -{Stealth[length=2.4mm]}},
  txc el/.style={bkBlue!75!black, line width=1.2pt, dash pattern=on 3pt off 1.5pt, -{Stealth[length=2.4mm]}},
  txc dim/.style={bkGray, line width=0.5pt, {Stealth[length=1.6mm]}-{Stealth[length=1.6mm]}},
  txc lvl/.style={bkInk, line width=1.1pt}}
% 電子（−の丸）  \txcel{(中心)}{半径}
\newcommand\txcel[2]{\filldraw[fill=bkBlue!70!black, draw=bkInk, line width=0.4pt] #1 circle (#2);\node[font=\sffamily\bfseries\scriptsize, text=white, inner sep=0pt] at #1 {$-$};}
\makeatother

\phchapter{2}{光はつぶでもある：\\光電効果と光子}{熱した鉄は、なぜ温度で光の色が変わる？}{\bklevelrange{2}{5}}{光はつぶでもある：光電効果と光子}
\bkmichishirube{5-7,5-10,5-14,7-10,熱-13,熱-15,波-6,波-13,電-2,電-15,原-1}{原-3,原-4,原-8,原-9,原-11,原-13}

\begin{goal}
\begin{itemize}[label={\color{bkTheme}●}, leftmargin=1.6\zw, itemsep=0pt, topsep=1pt]
\item \bklevel{2}\enspace 熱した物の色が温度で変わるわけを、光の強さの分布の山の動きで説明できる。
\item \bklevel{3}\enspace 光子のエネルギー $E=hf$ と運動量 $p=h/\lambda$ を求め、光電効果を $K_0=hf-W$ で説明できる。
\item \bklevel{3}\enspace コンプトン効果の式を読み、散乱した X 線・γ 線の波長とエネルギーを求められる。
\item \bklevel{5}\enspace 箱の中の光のモードを数えて、プランクの式・ウィーンの変位則・シュテファン・ボルツマン定数を導ける。
\end{itemize}
\end{goal}

\phkoushiki{%
\phk{光子のエネルギー $E=hf=hc/\lambda$}{高校物理}{ph:原-2:photon}
\phk{光電効果 $K_0=hf-W$、限界振動数 $f_0=W/h$}{高校物理}{ph:原-2:einstein}
\phk{阻止電圧 $eV_0=K_0$}{高校物理}{ph:原-2:stop}
\phk{光子の運動量 $p=h/\lambda=E/c$}{高校物理}{ph:原-2:momentum}
}

\begin{hatena}
\noindent\begin{minipage}[c]{0.56\linewidth}\raggedright
ガスコンロの上で鉄の棒を熱していくと、はじめは何も見えませんが、手を近づけると熱く感じます。もっと熱すると、暗い赤に光り始めます。さらに熱すると、明るいだいだい色、黄色、そして白っぽく光るようになります。

温度が上がると、なぜ光の色が変わっていくのでしょう。そして、なぜ、はじめは赤なのでしょう。青や緑から光り始めないのは、どうしてでしょう。
\end{minipage}\hfill
\begin{minipage}[c]{0.40\linewidth}\centering
\begin{tikzpicture}[x=1mm, y=1mm]
  \foreach \i/\c/\t in {0/bkRed!55!black/約\qty{900}{K}, 1/bkOrange!80!bkRed/約\qty{1300}{K}, 2/bkYellow!70!bkOrange/約\qty{1500}{K}} {
    \fill[\c] (0,{-\i*11}) rectangle (26,{-\i*11+5});
    \draw[bkInk, line width=0.4pt] (0,{-\i*11}) rectangle (26,{-\i*11+5});
    \node[txc lab, anchor=west] at (27.5,{-\i*11+2.5}) {\t};}
  \node[txc lab, anchor=south] at (22,6) {熱した鉄の棒の色（目安）};
\end{tikzpicture}
\end{minipage}
\end{hatena}

\noindent 原-1 では、電子と原子核が見つかった実験をたどりました。この章の主役は\textbf{光}です。物理シリーズの『波動』（波-13）と『電磁気学』（電-15）では、光は電場と磁場の波（電磁波）だと学びました。干渉も回折も、波だからこそ起こります。ところが、20世紀のはじめ、光が波であるだけでは説明できない現象が、2つ見つかりました。1つは、熱した物が出す光の色の分布（はてな？）、もう1つは、光を当てると金属から電子が飛び出す現象（\textbf{光電効果}）です。どちらも、「光は、決まった大きさのエネルギーの\textbf{つぶ}としてしか、出したり吸ったりできない」と考えると説明できます。

\begin{yattemiyou}[分光シートで、いろいろな明かりを見る]
\textbf{用意するもの}：分光シート（回折格子のシート。理科の教材）、白熱電球（なければ電球色の LED 電球）、昼光色の LED 電球、蛍光灯\par\vspace{2pt}
\begin{enumerate}[label={\bkcircled[bkGreen]{\arabic*}}, leftmargin=2.4\zw, itemsep=1pt, topsep=0pt]
\item 部屋を暗くして、分光シートを目の前に持ち、明かりから\qty{1}{m}ほど離れて見る。明かりの横に、色の帯（スペクトル）が見える。
\item 白熱電球・LED 電球・蛍光灯で、色の帯をくらべる。赤から紫まで、どの色が明るいかをノートにかく。
\item 白熱電球に調光器があれば、暗くしていく。色の帯の、どちらのはしから消えていくか。
\end{enumerate}
\par\vspace{2pt}\noindent\textbf{注意}：太陽は、分光シートを通しても、ぜったいに直接見ない。
\end{yattemiyou}

\noindent やってみようの答え：①②白熱電球は、赤から紫までとぎれのない帯で、赤い側が明るく、青・紫の側が暗く見えます。蛍光灯は、とぎれた明るい線（緑・だいだい・青など）が並びます（原-4）。LED 電球は、青の所に明るい山があり、黄色から赤にかけてなだらかな帯が続きます。③電球を暗くする（フィラメントの温度を下げる）と、\textbf{青・紫の側から先に消え}、最後に赤だけが残ります。温度が下がると、短い波長の光が先に出なくなるのです。そのわけを、この章で考えます。

\begin{omoidasou}
\bkrecall{熱-13}{黒体が出す光の分布は、プランクの式 $E_{\mathrm b\lambda}$ で表される（2つの定数 $c_1$・$c_2$ は測定から決めた）。山の波長はウィーンの変位則 $\lambda_{\max}T=\qty{2898}{\micro m.K}$、全体の強さは $\sigma T^4$。}
\bkrecall{熱-15}{エネルギーの粒 $\Delta\varepsilon$ の何個分かでしかエネルギーをもてないゆれは、温度 $T$ で、平均 $\Delta\varepsilon/(e^{\Delta\varepsilon/(k_{\mathrm B}T)}-1)$ のエネルギーをもつ（熱-15 ⑧）。}
\bkrecall{波-6}{長さ $L$ の弦の両はしを止めると、定在波の波長は $\lambda_n=2L/n$（$n=1,\ 2,\ 3,\ \ldots$）だけになる。}
\bkrecall{電-2}{電荷 $e$ の粒が電圧 $V$ で加速されると、運動エネルギーが $eV$ 変わる。$\qty{1}{eV}=\qty{1.60e-19}{J}$。}
\bkrecall{電-15}{光（電磁波）は、エネルギーだけでなく運動量も運ぶ。全部吸う面への放射圧は、光の強さ÷光速（$I/c$）。}
\bkrecall{7-10}{等比級数 $1+x+x^2+\cdots=1/(1-x)$（$|x|<1$。7-10⑤）。}
\end{omoidasou}

\noindent この章の記号：光の振動数は $f$ と書きます（物理シリーズの『波動』と同じ。高校の教科書は $\nu$（ニュー）と書く）。\textbf{仕事関数}を $W$ と書くので、この章では仕事を記号にしません。電子の運動エネルギーの最大値は $K_0$ と書きます。

%% ======================== 物理の道具 ========================
\bksection{doguu}{物理の道具}
\phstage{me}{光の色と、飛び出す電子を見る}{2}

\bksubsection{1}{熱した物の色と分光}
\label{ph:原-2:sec1}
熱した物は、温度で決まった分布の光を出します。光を全部吸う理想の物（\textbf{黒体}）が出す光の分布は、物の種類によらず、温度だけで決まります（熱-13）。図\ref{ph:原-2:fig-bunpu}は、温度ごとの分布を、山の高さが同じになるように縮めて重ねたものです。

\begin{bkfigure}
\begin{tikzpicture}
\begin{axis}[bkplot, width=112mm, height=56mm, xmin=0, xmax=4.2, ymin=0, ymax=1.45, xtick={0.5,1,1.5,2,2.5,3,3.5,4}, ytick={0.5,1}, xlabel={波長 $\lambda$（\unit{\micro m}）}, ylabel={強さ（山を1にそろえた）}]
  \fill[bkYellow!35] (axis cs:0.38,0) rectangle (axis cs:0.78,1.45);
  \addplot[bkRed!80!black, line width=1.5pt, domain=0.6:4.15, samples=120] {(2.898/x)^5*142.3/(exp(14.388/x)-1)};
  \addplot[bkOrange!85!black, line width=1.5pt, dash pattern=on 5pt off 2pt, domain=0.4:4.15, samples=140] {(1.932/x)^5*142.3/(exp(9.592/x)-1)};
  \addplot[bkGreen!50!black, line width=1.5pt, dash pattern=on 1.5pt off 1.5pt, domain=0.2:4.15, samples=160] {(1.035/x)^5*142.3/(exp(5.1386/x)-1)};
  \addplot[bkBlue!80!black, line width=1.5pt, dash pattern=on 6pt off 2pt on 1.5pt off 2pt, domain=0.12:4.15, samples=200] {(0.502/x)^5*142.3/(exp(2.4927/x)-1)};
  \node[font=\sffamily\footnotesize, anchor=north] at (axis cs:0.58,1.43) {可視光};
  \node[font=\sffamily\footnotesize, anchor=south, text=bkGreen!50!black] at (axis cs:1.2,1.04) {\qty{2800}{K}};
  \node[font=\sffamily\footnotesize, anchor=south, text=bkOrange!85!black] at (axis cs:1.93,1.04) {\qty{1500}{K}};
  \node[font=\sffamily\footnotesize, anchor=south, text=bkRed!80!black] at (axis cs:2.9,1.04) {\qty{1000}{K}};
  \node[font=\sffamily\footnotesize, anchor=south, text=bkBlue!80!black] at (axis cs:0.45,1.04) {\qty{5772}{K}};
\end{axis}
\end{tikzpicture}
\bkfigcaption{黒体の光の分布（山の高さをそろえて描いた）。温度が上がると山が短い波長へ動く。実線は \qty{1000}{K}、破線は \qty{1500}{K}、点線は \qty{2800}{K}（電球）、一点鎖線は \qty{5772}{K}（太陽）。黄色の帯が目に見える光（約 \qtyrange{0.38}{0.78}{\micro m}）。}\label{ph:原-2:fig-bunpu}
\end{bkfigure}

\noindent 山の波長は、ウィーンの変位則 $\lambda_{\max}=\qty{2898}{\micro m.K}/T$ で決まります。\qty{1000}{K}の鉄の山は \qty{2.9}{\micro m} の赤外線で、目に見える光（可視光）は、分布の左のすそのほんの少しです。そのすそでは、赤（長い波長）のほうが青（短い波長）よりずっと強いので、鉄は暗い赤に見えます（\qty{1000}{K}では、\qty{650}{nm}の赤の強さは \qty{450}{nm}の青の約3000倍）。温度が上がると、山が短い波長へ動いて、だいだい・黄色・緑・青の光も出るようになり、色は白に近づきます。太陽の表面（約 \qty{5772}{K}）では、山がちょうど可視光のまん中（約 \qty{0.50}{\micro m}）に来ます。

\begin{phmijika}[星の色と、太陽の表面の温度]
夜空の星にも、赤っぽい星と青白い星があります。赤い星（ベテルギウスなど）は表面の温度が低く（約 \qty{3500}{K}）、青白い星（リゲルなど）は高い（1万 K 以上）のです。星の光を分光して、どの波長がいちばん強いかを調べると、遠い星の表面の温度がわかります。熱-13 で学んだ太陽の表面の温度 \qty{5772}{K} も、この方法で求めたものです。暖房器具や電気コンロの熱線、トースターの赤い光も、同じしくみの色です。
\end{phmijika}

\noindent では、\textbf{なぜ}、温度が低いと短い波長の光が出ないのでしょう。電磁気学の考えでは、熱でゆれる物の中の電荷が、どんな波長の光でも出せるはずです。それどころか、③〜⑩で見るように、光を波だとだけ考えると、短い波長ほど多くのエネルギーが出ることになってしまい、実際とまったく合いません。この食いちがいを解いたのが、1900年のプランクです。答えは③と⑪で示します。

\bksubsection{2}{光を当てると、電子が飛び出す}
\label{ph:原-2:sec2}
光を金属に当てると、金属の表面から電子が飛び出すことがあります。これを\textbf{光電効果}といい、飛び出した電子を\textbf{光電子}といいます。\bkidx{こうでんこうか}{光電効果}1887年にヘルツが、紫外線を当てると放電が起こりやすくなることに気づいたのが、はじまりです。翌年、ハルヴァックスは、次のような実験で確かめました。

\begin{bkfigure}
\begin{tikzpicture}[x=1mm, y=1mm]
  \foreach \sx/\res/\dx in {0/はくが閉じる/0.8, 64/はくは開いたまま/4} {
  \begin{scope}[shift={(\sx,0)}]
  \draw[bkInk, line width=0.7pt] (0,0) rectangle (24,26);
  \draw[bkInk, line width=1pt] (12,26) -- (12,34);
  \fill[bkGray!60] (2,34) rectangle (22,36);
  \node[txc lab, anchor=west, align=left] at (24,35) {亜鉛の板\\（－に帯電）};
  \draw[bkInk, line width=0.8pt] (12,26) -- (12,12);
  \fill[bkYellow!70!black] (12,18) -- ({12+\dx},8) -- ({13+\dx},8.5) -- (12.6,18) -- cycle;
  \fill[bkYellow!70!black] (12,18) -- ({12-\dx},8) -- ({11-\dx},8.5) -- (11.4,18) -- cycle;
  \node[txc lab, anchor=north] at (12,-6) {\res};
  \end{scope}}
  \draw[txc ph] (-14,56) -- (6,38.5);
  \node[txc lab, anchor=south, text=bkOrange!85!black] at (-14,56.5) {紫外線};
  \foreach \p in {(9,40),(15,41.5)} {\txcel{\p}{1.3}}
  \draw[txc el] (15,43.5) -- (19,52);
  \draw[bkRed!75!black, line width=2.4pt, -{Stealth[length=3mm]}] (50,56) -- (70,38.5);
  \node[txc lab, anchor=south, text=bkRed!75!black] at (50,56.5) {強い赤い光};
  \node[txc lab, anchor=north, align=center] at (12,-1) {紫外線を当てると、すぐに};
  \node[txc lab, anchor=north, align=center] at (76,-1) {赤い光は、どんなに強くても};
\end{tikzpicture}
\bkfigcaption{はく検電器と亜鉛の板の実験（模式図。見るだけの実験）。紫外線が当たると、板から電子が飛び出して負の電荷が減り、はくが閉じる。}\label{ph:原-2:fig-kenden}
\end{bkfigure}

\noindent 調べていくと、光電効果には、光が波だと考えると、ふしぎなきまりがあることがわかってきました（④で詳しく見ます）。
\begin{enumerate}[label={\bkcircled[bkBlue]{\arabic*}}, leftmargin=2.4\zw, itemsep=1pt, topsep=2pt]
\item 金属ごとに決まった振動数より低い振動数の光（波長の長い光）では、どんなに強い光を当てても、電子は1つも飛び出さない。
\item 飛び出す電子の運動エネルギーの最大値は、光の強さによらず、光の振動数だけで決まる。
\item 光を強くすると、飛び出す電子の\textbf{数}が増える。
\item 光が当たると、どんなに弱い光でも、すぐに電子が飛び出す（待ち時間がない）。
\end{enumerate}
波なら、強い（振幅の大きい）光ほど、電子を大きくゆさぶり、勢いよく飛び出させるはずです。弱い光なら、エネルギーがたまるまで待たないと、電子は飛び出せないはずです。どちらも、実験と合いません。

\begin{column}{安全：紫外線の実験は、見るだけ}
紫外線のランプは、目や皮ふを傷めます。殺菌灯やブラックライトの光を、のぞきこんだり、長く当てたりしないでください。はく検電器と紫外線の実験は、先生が見せる実験です。レーザーポインターの光も、のぞいたり、人に向けたりしないでください（原-4・原-5）。
\end{column}

\Needspace{12\baselineskip}\phstage{te}{光子1個のエネルギーで考える}{3}

\bksubsection{3}{光子のエネルギー：プランク定数}
\label{ph:原-2:sec3}
1900年、ドイツのプランクは、熱した物が出す光の分布（①）を説明するために、大胆な仮定をおきました。「振動数 $f$ の光のエネルギーは、$hf$ という決まった大きさの\textbf{かたまり}の何個分かでしか、出たり入ったりしない」。$h$ は\textbf{プランク定数}とよばれる、とても小さな数です（1918年にノーベル物理学賞）。\bkidx{ぷらんくていすう}{プランク定数}
\[ h=6.626\,070\,15\times10^{-34}\ \mathrm{J\cdot s} \]
今では、この値をそのまま、質量の単位キログラムを決めるのに使っています（2019年から）。この本では $h=\qty{6.63e-34}{J.s}$ を使います。

1905年、アインシュタインは、さらに一歩進めて、「光そのものが、エネルギー $hf$ をもつつぶの集まりである」と考えました。このつぶを\textbf{光子}（光量子）といいます。\bkidx{こうし}{光子}光の速さ $c$、波長 $\lambda$ と振動数の関係 $c=f\lambda$（波-5）を使うと、光子1個のエネルギーは
\pheq{E=hf=\frac{hc}{\lambda}}{J ＝ J·s × Hz}
\label{ph:原-2:photon}です。高校の教科書では、振動数を $\nu$（ニュー）と書いて $E=h\nu$ とします。同じ式です。

\begin{phyomi}[$E=hf$]
「光は、振動数に比例した大きさのつぶ」。振動数の高い（波長の短い）光ほど、1つぶのエネルギーが大きい。光の強さ（明るさ）は、つぶの\textbf{数}で決まります。同じ明るさでも、青い光はつぶが大きく、数が少ない。赤い光はつぶが小さく、数が多いのです。
\end{phyomi}

\phboxhead{電子ボルトで書くと便利}
光子のエネルギーは、J で書くととても小さな数になります。そこで、電子ボルト（電-2）を使います。$hc$ を eV と nm で書くと、
\[ hc=\frac{6.63\times10^{-34}\times3.00\times10^8}{1.60\times10^{-19}}\ \mathrm{eV\cdot m}\approx\qty{1.24e-6}{eV.m}=\qty{1240}{eV.nm} \]
なので、波長を nm で入れれば $E\,[\mathrm{eV}]=1240/\lambda\,[\mathrm{nm}]$ です（例題2）。可視光の光子は、赤（\qty{780}{nm}）の \qty{1.6}{eV} から紫（\qty{380}{nm}）の \qty{3.3}{eV} までです。

\begin{phbaisuru}{光の振動数 $f$}
光子1個のエネルギーも2倍になります。$E$ は $f$ に比例し、波長 $\lambda$ には反比例します。下の図は、光の色と光子1個のエネルギーです。
\par\smallskip\noindent\begin{tikzpicture}[x=1mm, y=1mm]
  \foreach \l/\c in {380/bkPurple!80, 420/bkPurple!60!bkBlue, 460/bkBlue!80, 500/bkGreen!60!bkBlue, 540/bkGreen!70, 580/bkYellow!80!bkOrange, 620/bkOrange!90, 660/bkRed!80, 700/bkRed!65!black, 740/bkRed!45!black} {
    \fill[\c] ({(\l-380)*0.2},0) rectangle ({(\l-380)*0.2+8},5);}
  \draw[bkInk, line width=0.4pt] (0,0) rectangle (80,5);
  \foreach \l in {400,500,600,700} {\draw[bkInk, line width=0.4pt] ({(\l-380)*0.2},0) -- ({(\l-380)*0.2},-1.2); \node[txc lab, anchor=north] at ({(\l-380)*0.2},-1.4) {\l};}
  \node[txc lab, anchor=west] at (82,-3.3) {波長（nm）};
  \foreach \l/\e in {400/3.10,500/2.48,600/2.07,700/1.77} {\node[txc lab, anchor=south] at ({(\l-380)*0.2},6) {\e};}
  \node[txc lab, anchor=west] at (82,8) {光子1個（eV）};
\end{tikzpicture}
\end{phbaisuru}

\begin{phketa}[光子のエネルギーの目安]
\phketaline{FM ラジオ（\qty{80}{MHz}）の光子}{約 $3\times10^{-7}$\,eV}
\phketaline{室温の熱のエネルギー $k_{\mathrm B}T$（\qty{300}{K}）}{\qty{0.0259}{eV}}
\phketaline{赤外線（\qty{10}{\micro m}）・可視光}{\qty{0.12}{eV}・\qtyrange{1.6}{3.3}{eV}}
\phketaline{X 線（\qty{0.1}{nm}）・γ 線}{約 \qty{12}{keV}・約 \qty{1}{MeV}}
\end{phketa}

\bksubsection{4}{光電効果のきまりを、光子で説明する}
\label{ph:原-2:sec4}
光が光子のつぶだと考えると、②の4つのきまりは、すっきり説明できます。かぎは、「\textbf{光子1個のエネルギーは、電子1個にまとめて渡される}」ことです。

金属の中の電子は、金属の中を自由に動けますが、表面から外へ出るには、エネルギーが要ります。外へ出るのに要る、いちばん小さいエネルギーを、その金属の\textbf{仕事関数} $W$ といいます。\bkidx{しごとかんすう}{仕事関数}金属の中の電子は、深さ $W$ の「くぼみ」の底にいるようなものです。

\begin{bkfigure}
\begin{tikzpicture}[x=1mm, y=1mm]
  \fill[bkGray!20] (0,0) rectangle (40,30);
  \draw[txc lvl] (0,30) -- (40,30) -- (40,46) -- (100,46);
  \node[txc lab, anchor=north west, align=left] at (62,44.5) {金属の外\\（止まった電子のエネルギー）};
  \node[txc lab, anchor=north] at (20,-1) {金属};
  \foreach \x in {6,14,22,30} {\txcel{(\x,30)}{1.5}}
  \draw[txc dim] (36,30.5) -- (36,45.5); \node[txc lab, anchor=east] at (35,38) {$W$};
  \draw[txc ph] (-24,14) -- (20,28);
  \node[txc lab, anchor=north, text=bkOrange!85!black] at (-14,12) {光子 $hf$};
  \draw[bkPurple!80!black, line width=1.2pt, -{Stealth[length=2.4mm]}] (22,31.5) -- (22,58);
  \txcel{(22,60)}{1.5}
  \draw[txc dim] (50,46.5) -- (50,60); \node[txc lab, anchor=west, align=left] at (51,53) {$K_0=hf-W$\\（外での運動エネルギー）};
  \draw[bkGray, line width=0.4pt, dash pattern=on 1.5pt off 1.5pt] (23.5,60) -- (50,60);
  \draw[txc dim] (14,30.5) -- (14,59.5); \node[txc lab, anchor=east] at (13,52) {$hf$};
\end{tikzpicture}
\bkfigcaption{光電効果のエネルギーの図。光子1個のエネルギー $hf$ を受けとった電子は、仕事関数 $W$ の分を使って外へ出て、残り $hf-W$ を運動エネルギーとしてもつ（表面にいちばん近い、浅い所の電子のとき）。}\label{ph:原-2:fig-energy}
\end{bkfigure}

\begin{enumerate}[label={\bkcircled[bkBlue]{\arabic*}}, leftmargin=2.4\zw, itemsep=1pt, topsep=2pt]
\item $hf<W$ の光子は、1個では電子を外へ出せません。光を強くしても、光子の数が増えるだけで、1個1個の光子のエネルギーは小さいままなので、電子は出てきません。
\item 外へ出た電子の運動エネルギーは、光子1個のエネルギーから $W$ を引いた残りで、光子の数（光の強さ）にはよりません。
\item 光を強くすると光子の数が増え、光子を受けとる電子の数が増えます。
\item 光子1個のエネルギーは、一度に電子に渡されるので、待ち時間はありません。
\end{enumerate}

\begin{bkcenter}
\begin{bktable}{colspec={Q[l,m,wd=34mm] *{4}{Q[c,m,wd=21mm]}}, rowsep=1pt, row{1}={bg=bkTheme!10}}
金属 & セシウム & カリウム & タングステン & 銅 \\
仕事関数 $W$（目安） & 約\qty{2.1}{eV} & 約\qty{2.3}{eV} & 約\qty{4.5}{eV} & 約\qty{4.7}{eV} \\
限界の波長 $hc/W$ & 約\qty{590}{nm} & 約\qty{540}{nm} & 約\qty{280}{nm} & 約\qty{260}{nm} \\
\end{bktable}
\bktabcaption{仕事関数の目安。同じ金属でも、表面の向きや汚れで 0.1〜0.5\,eV ほど違う。セシウムやカリウムは可視光でも電子が飛び出すが、銅やタングステンは紫外線でないと飛び出さない。}\label{ph:原-2:tab-W}
\end{bkcenter}

\bksubsection{5}{アインシュタインの式と阻止電圧}
\label{ph:原-2:sec5}
図\ref{ph:原-2:fig-energy}のエネルギーの保存を式にすると、飛び出す電子の運動エネルギーの最大値 $K_0$ は
\pheq{K_0=hf-W}{J ＝ J·s × Hz － J}
\label{ph:原-2:einstein}です（\textbf{アインシュタインの光電効果の式}。1921年にノーベル物理学賞）。「最大値」というのは、金属の深い所にいる電子は、外へ出るまでにもっとエネルギーを使うので、それより遅く出てくるからです。$K_0=0$ になる振動数を\textbf{限界振動数} $f_0$ といいます。
\[ hf_0=W\qquad\text{より}\qquad f_0=\frac{W}{h},\qquad\text{限界の波長}\quad\lambda_0=\frac{c}{f_0}=\frac{hc}{W} \]
$f<f_0$（$\lambda>\lambda_0$）の光では、電子は飛び出しません。

\begin{bkfigure}
\begin{tikzpicture}[x=1mm, y=1mm]
  \draw[bkInk, line width=0.8pt, rounded corners=6mm] (0,0) rectangle (56,34);
  \draw[bkInk, line width=1.6pt] (8,6) arc (-35:35:19.2);
  \node[txc lab, anchor=north] at (10,-1) {光電面};
  \draw[bkInk, line width=1.6pt] (46,8) -- (46,26);
  \node[txc lab, anchor=north] at (46,-1) {電極};
  \draw[txc ph] (28,48) -- (14,24);
  \node[txc lab, anchor=south, text=bkOrange!85!black] at (30,48) {光（振動数 $f$）};
  \draw[txc el] (16,14) -- (40,14);
  \node[txc lab, anchor=north] at (28,12.5) {電子};
  \draw[bkInk, line width=0.7pt] (11.6,17) -- (-10,17) -- (-10,-14) -- (12,-14);
  \draw[bkInk, line width=0.7pt] (15,-14) -- (28,-14);
  \draw[bkInk, line width=0.7pt] (36,-14) -- (66,-14) -- (66,17) -- (46,17);
  \draw[bkInk, line width=0.7pt] (32,-14) circle (4); \node[txc lab] at (32,-14) {A};
  \draw[bkInk, line width=1.6pt] (12,-9) -- (12,-19); \draw[bkInk, line width=0.9pt] (15,-11) -- (15,-17);
  \draw[bkInk, line width=0.6pt, -{Stealth[length=2mm]}] (7,-21) -- (20,-7);
  \node[txc lab, anchor=north, align=center] at (13,-22) {電圧 $V$ を変えられる電源};
  \node[txc lab, anchor=west, align=left] at (70,8) {電極を光電面より\\低い電位にすると、\\電子は押しもどされる};
\end{tikzpicture}
\bkfigcaption{光電管で阻止電圧をはかる回路（模式図）。A は電流計。電極の電位を光電面より $V$ だけ低くし、$V$ を大きくしていくと、ある電圧 $V_0$ で電流が0になる。}\label{ph:原-2:fig-kokan}
\end{bkfigure}

\noindent $K_0$ は、図\ref{ph:原-2:fig-kokan}の光電管ではかれます。電極の電位を、光電面より $V$ だけ低くすると、電子は電極へ向かう途中で、電場に押しもどされ、運動エネルギーを $eV$ 失います（電-2）。いちばん速い電子でも届かなくなる電圧 $V_0$（\textbf{阻止電圧}）では、
\pheq{eV_0=K_0\qquad\text{より}\qquad V_0=\frac{h}{e}f-\frac{W}{e}}{V ＝ (J·s/C) × Hz － J/C}
\label{ph:原-2:stop}です。\bkidx{そしでんあつ}{阻止電圧}$V_0$ を V で書いた数と、$K_0$ を eV で書いた数は同じです（$K_0=\qty{0.80}{eV}$ なら $V_0=\qty{0.80}{V}$）。

\begin{bkfigure}
\begin{tikzpicture}
\begin{axis}[bkplot, width=62mm, height=46mm, xmin=-2.4, xmax=3.2, ymin=0, ymax=2.4, xtick={-2,0,1,2,3}, ytick=\empty, xlabel={電圧（電極の電位－光電面）（V）}, ylabel={電流}, axis lines=left]
  \addplot[bkcurve, domain=-0.8:3.1, samples=80] {2*(1-exp(-(x+0.8)/0.45))};
  \addplot[bkOrange!85!black, line width=1.5pt, dash pattern=on 5pt off 2pt, domain=-0.8:3.1, samples=80] {1*(1-exp(-(x+0.8)/0.45))};
  \addplot[bkGreen!50!black, line width=1.5pt, dash pattern=on 1.5pt off 1.5pt, domain=-1.9:3.1, samples=80] {1*(1-exp(-(x+1.9)/0.55))};
  \node[font=\sffamily\footnotesize, anchor=north, inner sep=2pt] at (axis cs:-0.8,0) {$-V_0$};
  \node[font=\sffamily\footnotesize, anchor=south west, text=bkBlue!80!black] at (axis cs:0.6,2.02) {強い光};
  \node[font=\sffamily\footnotesize, anchor=south west, text=bkOrange!85!black] at (axis cs:1.0,1.02) {弱い光};
\end{axis}
\end{tikzpicture}\hspace{3mm}
\begin{minipage}[b]{52mm}\small\raggedright
実線（強い光）と破線（弱い光）は同じ振動数の光。電流の大きさ（電子の数）は2倍違うが、電流が0になる電圧 $-V_0$ は同じ。\par\vspace{3pt}
点線は、振動数の高い光（弱い光）。$V_0$ が大きくなる（飛び出す電子が速い）。\par\vspace{6pt}
\end{minipage}
\bkfigcaption{光電管の電流と電圧（形の模式図）。阻止電圧は光の強さによらず、振動数で決まる。}\label{ph:原-2:fig-iv}
\end{bkfigure}

\begin{phflow}[エネルギーの保存から光電効果の式へ]
\phflowrow{$hf=W+K$}{いちばん浅い所の\\電子（$K$ が最大）}{$K_0=hf-W$}
\phflowrow{$K_0=hf-W$}{ちょうど\\飛び出せない}{$f_0=\dfrac{W}{h}$}
\phflowrow{$K_0=hf-W$}{電子を電圧で\\止める}{$eV_0=K_0$}
\end{phflow}

\begin{phyomi}[$K_0=hf-W$]
「光子1個のエネルギーから、金属を出る『入場料』$W$ を引いた残りが、電子の勢い」。お金にたとえると、光子は決まった額のお札、$W$ は出口の料金です。料金より安いお札は、何枚あっても1枚ずつしか使えないので、出口を通れません。
\end{phyomi}

\begin{phbaisuru}{光の強さ}
飛び出す電子の数（電流）は2倍になりますが、$K_0$ と $V_0$ は変わりません（図\ref{ph:原-2:fig-iv}）。一方、\textbf{振動数}を2倍にすると、光子のエネルギー $hf$ は2倍になり、$K_0=hf-W$ は2倍よりも\textbf{大きく}増えます（$W$ の分を引く前の量が2倍になるから）。
\end{phbaisuru}

\Needspace{8\baselineskip}\phboxhead{ミリカンの測定：グラフの傾きが $h/e$}
\label{ph:原-2:millikan}
$V_0=(h/e)f-W/e$ は、横に $f$、たてに $V_0$ をとると直線です。1916年、ミリカンは、いろいろな振動数の光で $V_0$ をはかり、直線になることを確かめました。直線の\textbf{傾き}は $h/e$ なので、傾きに $e$ をかけると $h$ が求まります。ミリカンの値は、プランクが熱した物の光から求めた $h$ と、よく合いました。直線を横の軸まで延ばした所が限界振動数 $f_0$、たての軸との交わる所が $-W/e$ です。傾きは、金属の種類によらず、どれも同じ $h/e$ です（例題5）。

\begin{bkfigure}
\begin{tikzpicture}
\begin{axis}[bkplot, width=98mm, height=54mm, xmin=0, xmax=10.2, ymin=-2.6, ymax=2.1, xtick={2,4,6,8,10}, xticklabels={2,,6,8,10}, ytick={-2,-1,0,1,2}, xlabel={光の振動数 $f$（$10^{14}$\,Hz）}, ylabel={阻止電圧 $V_0$（V）}]
  \draw[bkInk, line width=0.6pt] (axis cs:0,0) -- (axis cs:10.2,0);
  \addplot[bkcurve, domain=4.85:10, samples=2] {0.4133*x-2.003};
  \addplot[bkcurve, dashed, line width=0.9pt, domain=0:4.85, samples=2] {0.4133*x-2.003};
  \addplot[bkOrange!85!black, line width=1.5pt, dash pattern=on 5pt off 2pt, domain=5.56:10, samples=2] {0.4136*x-2.3};
  \fill[bkRed] (axis cs:5.5,0.27) circle (2pt) (axis cs:6.5,0.69) circle (2pt) (axis cs:7.5,1.10) circle (2pt) (axis cs:8.5,1.51) circle (2pt);
  \node[font=\sffamily\footnotesize, anchor=south east] at (axis cs:4.8,0.08) {$f_0$};
  \node[font=\sffamily\footnotesize, anchor=north west] at (axis cs:0.15,-2.15) {$-W/e$};
  \node[font=\sffamily\footnotesize, anchor=west, text=bkBlue!80!black] at (axis cs:0.3,1.65) {実線：例題5の金属};
  \node[font=\sffamily\footnotesize, anchor=west, text=bkOrange!85!black] at (axis cs:0.3,1.1) {破線：カリウム};
\end{axis}
\end{tikzpicture}
\bkfigcaption{阻止電圧と光の振動数。実線が例題5の金属（$W=\qty{2.0}{eV}$。点は、はかった値）、破線がカリウム（$W\approx\qty{2.3}{eV}$）。2本は平行（傾きはどちらも $h/e$）で、ずれは仕事関数の違い。}\label{ph:原-2:fig-millikan}
\end{bkfigure}

\begin{phmodel}[光電効果の光子と電子]
\phignore{光電面の表面の汚れや、電極の金属の仕事関数の違い（電極と光電面の間の小さな電圧）。}
\phassume{光子1個のエネルギーは、電子1個に全部渡される。1個の電子は、光子を1個だけ吸う。}
\phlimit{とても強いレーザーの光で、1個の電子が同時に2個以上の光子を吸うとき（$hf<W$ でも電子が出ることがある）。}
\end{phmodel}

\bksubsection{6}{光子の運動量}
\label{ph:原-2:sec6}
『電磁気学』（電-15 ⑤）で、光は運動量も運び、エネルギー $U$ の光は運動量 $U/c$ をもつことを学びました。光子1個のエネルギーは $hf$ なので、光子1個の運動量は
\pheq{p=\frac{E}{c}=\frac{hf}{c}=\frac{h}{\lambda}}{kg·m/s ＝ J·s ÷ m}
\label{ph:原-2:momentum}です。光子には質量がないので、運動量は $mv$ では書けません。それでも運動量をもち、ぶつかった物を押します。強さ $I$（\qty{1}{m^2} に1秒あたり $I$ J）の光を全部吸う面には、1秒に $I/(hf)$ 個の光子が当たり、1個あたり $h/\lambda$ の運動量を渡すので、押す力は \qty{1}{m^2} あたり
\[ \frac{I}{hf}\times\frac{h}{\lambda}=\frac{I}{f\lambda}=\frac{I}{c} \]
となり、電-15 の放射圧と同じ答えになります（例題7）。

\begin{phbaisuru}[半分]{波長 $\lambda$}
光子1個の運動量は2倍になります。X 線（\qty{0.1}{nm}）の光子の運動量は、可視光（\qty{500}{nm}）の5000倍です。⑦のコンプトン効果のように、X 線の光子が電子をはじき飛ばせるのは、運動量が大きいからです。
\end{phbaisuru}

\bksubsection{7}{コンプトン効果：光子と電子のぶつかり}
\label{ph:原-2:sec7}
1923年、アメリカのコンプトンは、X 線を炭素（黒鉛）に当てると、散乱した X 線の中に、もとより\textbf{波長の長い} X 線がまじっていることを見つけました（1927年にノーベル物理学賞）。散乱した角が大きいほど、波長ののびは大きくなりました。波なら、散乱しても振動数（波長）は変わらないはずです。

コンプトンは、これを「光子が、止まっている電子とビリヤードの球のようにぶつかり、電子をはじき飛ばしてエネルギーを失う」と考えて説明しました。エネルギーが減ると、$E=hc/\lambda$ なので波長がのびます。エネルギーと運動量の保存（力-3）を、光速に近い電子にも使えるように相対論の形で書くと（原-11 ⑧で導く）、散乱角 $\theta$ の光子の波長ののびは
\pheq{\Delta\lambda=\lambda'-\lambda=\frac{h}{m_{\mathrm e}c}(1-\cos\theta)}{m ＝ J·s ÷ (kg × m/s)}
\label{ph:原-2:compton}になります。\bkidx{こんぷとんこうか}{コンプトン効果}$h/(m_{\mathrm e}c)=\qty{2.43}{pm}$（ピコメートル。$\qty{1}{pm}=10^{-12}$\,m）を、電子の\textbf{コンプトン波長} $\lambda_{\mathrm C}$ といいます。

\begin{bkfigure}
\begin{tikzpicture}[x=1mm, y=1mm]
  \draw[txc ph] (-46,0) -- (-3,0);
  \node[txc lab, anchor=south] at (-26,2.5) {X 線の光子（波長 $\lambda$）};
  \txcel{(0,0)}{2}
  \node[txc lab, anchor=north east] at (-1.5,-2.5) {止まっていた電子};
  \draw[txc ph, decoration={snake, amplitude=1.1mm, segment length=4.4mm, post length=2.4mm}] (2,2) -- (36,26);
  \node[txc lab, anchor=west] at (37,27) {散乱した光子（波長 $\lambda'$、$\lambda'>\lambda$）};
  \draw[bkBlue!75!black, line width=1.3pt, -{Stealth[length=2.4mm]}] (2.5,-2) -- (30,-18);
  \node[txc lab, anchor=west] at (31,-18) {はじき飛ばされた電子};
  \draw[bkGray, line width=0.5pt, dash pattern=on 2pt off 1.5pt] (2,0) -- (40,0);
  \draw[bkInk, line width=0.5pt] (14,0) arc (0:35.2:14);
  \node[txc lab, anchor=west] at (15,5) {$\theta$};
\end{tikzpicture}
\bkfigcaption{コンプトン効果（模式図）。光子は、電子にエネルギーと運動量を分けたぶん、エネルギーが減って波長がのびる。}\label{ph:原-2:fig-compton}
\end{bkfigure}

\begin{phyomi}[$\Delta\lambda=\lambda_{\mathrm C}(1-\cos\theta)$]
「大きく曲がった光子ほど、電子に多く渡して、波長が多くのびる」。のびは、もとの波長によらず、角だけで決まり、真横（$90^\circ$）で $\lambda_{\mathrm C}$、真後ろ（$180^\circ$）で $2\lambda_{\mathrm C}$ です。のびはピコメートルの大きさなので、可視光（\qty{500}{nm}）ではわずか $10^{-5}$ ほどで目立ちません。X 線やγ線のように波長の短い光でこそ、はっきり見えます。
\end{phyomi}

\begin{phbaisuru}{角の $1-\cos\theta$}
波長ののびは2倍になります。$1-\cos\theta$ は、$\theta=0^\circ$ で0、$90^\circ$ で1、$180^\circ$ で2です。
\par\smallskip\noindent\begin{tikzpicture}
\begin{axis}[bkplot, width=92mm, height=38mm, xmin=0, xmax=185, ymin=0, ymax=5.5, xtick={0,45,90,135,180}, ytick={1,2.43,4.85}, yticklabels={1,2.43,4.85}, xlabel={散乱角 $\theta$（度）}, ylabel={$\Delta\lambda$（pm）}]
  \addplot[bkcurve, domain=0:180, samples=90] {2.426*(1-cos(x))};
  \fill[bkRed] (axis cs:90,2.426) circle (1.8pt) (axis cs:180,4.853) circle (1.8pt);
\end{axis}
\end{tikzpicture}
\end{phbaisuru}

\noindent γ 線のように光子のエネルギーで考えるときは、$E=hc/\lambda$ を使って書きかえた
\[ E'=\frac{E}{1+\dfrac{E}{m_{\mathrm e}c^2}(1-\cos\theta)},\qquad m_{\mathrm e}c^2=\qty{0.511}{MeV} \]
が便利です（例題10）。はじき飛ばされた電子は、差の $E-E'$ のエネルギーをもちます。γ 線が物を通りぬけるときに起こるこのぶつかりは、原-13 で放射線と物のふるまいとして使います。

\bksubsection{8}{光子を1つずつ数える干渉じま}
\label{ph:原-2:sec8}
光が光子のつぶなら、波の性質である干渉はどうなるのでしょう。『波動』（波-13）の二重スリットの実験を、とても弱い光で行います。光子が1つずつしか通らないほど弱くして、光子を1個ずつ数えられるカメラでスクリーンを写すと、光子は1個ずつ、ばらばらの場所に点として届きます。ところが、点がたくさん集まると、だんだん干渉じまが見えてくるのです（図\ref{ph:原-2:fig-ichiko}）。1909年にテイラーは、写真の板を何か月も光に当てる方法で、とても弱い光でも干渉じまができることを確かめました。今では、光子を1個ずつ出す光源と、1個ずつ数える検出器で、この実験が確かめられています。

\begin{bkfigure}
\begin{tikzpicture}[x=1mm, y=1mm]
  \begin{scope}[shift={(0,0)}]
  \fill[bkInk!8] (0,0) rectangle (36,18); \draw[bkInk, line width=0.5pt] (0,0) rectangle (36,18);
  \foreach \p in {(16.9,9.7),(20.7,4.2),(28.7,11.0),(4.6,15.2),(7.5,17.1),(24.4,16.4),(13.0,3.0),(21.7,12.0),(12.2,14.4),(17.3,8.7),(25.4,17.0),(30.4,13.8),(13.2,0.8),(12.8,13.6),(19.6,5.9),(6.4,3.1),(5.8,17.2),(19.2,4.6),(16.4,1.5),(17.8,2.8)} {\fill[bkBlue!80!black] \p circle (0.45);}
  \node[txc lab, anchor=north] at (18,-1) {光子 20 個};
  \end{scope}
  \begin{scope}[shift={(43,0)}]
  \fill[bkInk!8] (0,0) rectangle (36,18); \draw[bkInk, line width=0.5pt] (0,0) rectangle (36,18);
  \foreach \p in {(17.6,8.6),(12.2,5.2),(23.9,6.7),(25.4,8.8),(23.4,0.8),(19.7,11.9),(12.5,13.0),(29.8,14.8),(18.7,3.4),(17.2,12.7),(24.0,11.1),(25.9,3.3),(22.7,8.5),(8.1,2.8),(11.4,3.8),(32.5,7.4),(13.5,10.9),(12.2,16.7),(29.5,14.3),(23.1,3.0),(32.2,2.4),(17.9,6.0),(23.2,13.1),(12.1,9.5),(16.7,13.1),(19.6,13.7),(30.7,5.6),(29.9,14.6),(25.0,5.4),(11.9,10.4),(29.2,4.5),(23.9,12.7),(25.2,14.7),(13.7,1.6),(31.6,7.9),(11.3,8.8),(18.2,9.9),(22.4,9.9),(18.6,4.5),(23.6,11.8),(8.0,16.9),(12.8,6.2),(13.9,14.6),(18.6,10.2),(7.3,15.5),(17.3,0.6),(16.8,3.2),(5.7,7.1),(5.8,4.6),(17.4,16.1),(7.2,4.5),(11.2,5.4),(5.9,11.9),(0.2,13.7),(18.7,2.2),(17.7,1.5),(14.3,4.9),(1.8,16.9),(30.8,0.9),(19.4,10.2),(11.1,14.5),(8.1,5.8),(29.5,2.8),(23.5,10.3),(13.8,1.7),(6.1,12.0),(10.3,13.4),(20.1,1.5),(22.0,7.8),(13.0,0.8),(26.2,9.5),(28.6,5.9),(29.3,1.5),(35.2,10.3),(17.4,7.2),(17.6,3.4),(13.6,9.2),(25.3,17.1),(13.1,12.0),(18.8,1.9),(25.8,13.3),(22.0,9.3),(14.2,7.2),(8.1,10.1),(1.1,10.1),(23.3,12.8),(5.6,3.2),(22.2,13.3),(19.6,2.6),(24.6,16.8),(10.7,8.6),(18.7,7.7),(29.2,13.5),(13.0,4.8),(16.5,9.3),(17.0,5.3),(18.3,7.2),(23.2,1.2),(16.9,11.8),(16.5,7.0),(7.1,14.5),(4.9,16.1),(33.8,3.7),(25.3,2.5),(23.2,6.8),(17.2,16.8),(17.3,0.7),(18.1,13.9),(12.4,5.6),(29.9,6.5),(17.8,4.1),(16.3,2.1),(12.3,16.9),(22.5,11.2),(4.9,13.0),(20.2,1.3),(31.4,10.5),(4.0,12.1),(30.4,13.3),(18.6,8.4),(5.6,5.0),(6.3,8.9),(24.7,10.2),(30.1,2.4),(31.3,5.7),(18.2,7.6),(8.1,10.1),(10.4,9.2),(16.8,13.5),(19.5,5.2),(18.5,7.0),(19.0,1.5),(4.8,12.3),(35.9,3.4),(12.1,11.3),(17.6,16.0),(31.6,7.7),(23.8,14.7),(24.3,12.9),(18.2,10.9),(28.7,8.9),(17.2,4.0),(23.4,11.4),(18.1,5.7),(30.3,1.6),(17.2,1.0),(14.0,16.6),(6.9,4.8),(30.2,2.2),(5.2,0.7),(3.8,9.8),(0.2,10.7),(19.3,7.8),(29.4,9.6),(24.4,9.3),(16.9,3.1),(12.3,12.1),(16.6,15.1),(11.8,8.2),(31.4,4.4),(10.3,14.4),(18.1,9.1),(11.0,10.7),(22.6,11.6),(5.8,12.3),(18.4,14.7),(11.5,2.5),(10.0,7.5),(16.4,1.2),(11.0,10.9),(1.2,3.8),(9.7,3.5),(10.2,5.0),(16.5,2.0),(5.4,10.4),(8.2,5.7),(23.9,5.3),(24.9,3.7),(16.8,2.2),(17.1,14.7),(11.2,8.5),(10.6,5.1),(5.8,15.9),(17.3,3.1),(24.0,8.6),(23.8,12.1),(22.3,13.3),(18.9,13.5),(23.5,13.9),(31.0,8.8),(29.8,7.8),(24.0,15.1),(17.8,11.6),(20.2,16.7),(11.5,1.3),(29.2,11.8),(17.2,5.0),(13.1,2.6),(5.4,13.5),(22.6,2.6)} {\fill[bkBlue!80!black] \p circle (0.35);}
  \node[txc lab, anchor=north] at (18,-1) {光子 200 個};
  \end{scope}
  \begin{scope}[shift={(86,0)}]
  \fill[bkInk!8] (0,0) rectangle (36,18); \draw[bkInk, line width=0.5pt] (0,0) rectangle (36,18);
  \foreach \p in {(18.9,3.0),(28.3,8.0),(34.1,15.6),(15.3,5.0),(11.0,7.8),(18.7,5.8),(18.9,11.9),(18.8,3.1),(1.1,7.4),(12.1,5.9),(17.7,14.9),(17.4,2.0),(7.5,1.7),(18.8,8.1),(17.7,2.9),(17.4,5.6),(24.5,6.4),(12.4,4.4),(18.6,5.1),(31.7,3.9),(7.0,2.5),(16.5,11.3),(12.9,5.7),(17.6,8.6),(5.6,12.1),(11.8,1.0),(12.9,12.0),(23.2,15.5),(17.0,11.6),(13.9,13.4),(25.1,0.6),(10.9,9.8),(23.4,7.5),(13.9,3.5),(23.6,14.3),(11.8,13.2),(24.5,15.5),(11.8,11.4),(14.0,12.1),(5.3,16.6),(17.9,6.4),(24.5,12.9),(5.4,11.2),(1.7,6.9),(24.3,10.4),(25.2,2.5),(16.9,6.7),(24.6,10.6),(7.3,9.8),(25.0,7.8),(6.0,6.2),(18.8,4.9),(5.1,8.0),(16.7,4.9),(15.4,5.1),(24.7,17.4),(17.8,9.9),(23.9,16.5),(34.2,1.8),(4.6,13.1),(7.0,11.0),(12.4,3.3),(14.1,11.6),(17.9,8.0),(11.8,11.2),(17.0,11.1),(19.4,11.8),(24.5,8.9),(16.2,3.6),(34.5,3.0),(25.4,2.1),(23.4,12.4),(18.0,3.0),(11.7,16.8),(28.9,11.9),(29.5,15.5),(22.5,8.4),(14.1,14.1),(11.8,9.8),(24.2,13.7),(10.7,6.1),(30.6,5.7),(17.9,4.1),(5.5,16.7),(11.1,15.2),(6.2,0.9),(12.2,8.2),(23.3,6.7),(7.2,13.2),(29.3,7.1),(17.9,2.7),(19.2,6.1),(17.7,11.2),(25.9,8.8),(14.6,3.6),(23.2,6.6),(23.7,13.7),(24.6,12.2),(23.1,12.3),(12.4,5.5),(11.9,1.0),(35.1,15.6),(6.4,16.5),(11.3,4.7),(22.4,8.7),(11.7,10.9),(18.3,9.1),(13.2,12.3),(18.1,10.4),(12.5,3.9),(17.5,0.9),(5.4,13.9),(6.8,2.9),(25.3,10.3),(19.1,9.3),(17.3,14.6),(12.0,11.3),(17.0,3.9),(18.7,8.4),(12.3,13.6),(24.7,16.3),(12.0,2.5),(30.1,5.8),(12.2,9.5),(35.5,2.9),(19.4,15.6),(29.7,8.4),(18.3,2.8),(13.2,2.0),(4.8,17.1),(6.0,4.4),(7.5,8.4),(27.9,8.9),(19.2,12.7),(17.8,8.9),(25.2,10.9),(6.9,1.8),(12.1,1.6),(18.4,17.1),(22.3,14.8),(17.2,2.7),(28.8,6.8),(29.7,3.7),(18.3,14.8),(3.7,6.5),(25.8,0.8),(0.3,16.6),(16.4,1.4),(23.1,3.6),(12.3,3.8),(20.1,16.1),(17.1,13.2),(12.3,4.1),(10.4,7.0),(30.1,12.9),(18.8,17.1),(18.2,15.3),(18.2,10.9),(12.5,17.0),(18.3,14.5),(11.3,16.4),(0.7,15.3),(17.8,17.3),(24.8,13.1),(24.1,13.8),(22.7,15.3),(22.2,9.7),(3.8,16.7),(34.7,14.6),(31.8,11.4),(24.2,4.6),(8.0,9.0),(19.2,13.2),(30.0,8.2),(18.3,16.6),(12.6,0.8),(25.5,15.7),(35.4,1.8),(23.8,3.0),(18.3,4.0),(24.8,8.2),(30.5,12.8),(16.8,2.1),(6.8,13.8),(19.3,6.4),(4.4,15.1),(31.0,1.9),(18.2,13.7),(28.6,14.3),(22.4,2.1),(10.9,3.7),(23.5,10.1),(25.2,15.6),(11.6,4.0),(12.2,10.2),(7.7,15.3),(17.2,16.3),(13.6,6.6),(18.5,16.3),(0.6,15.6),(23.0,16.4),(19.3,17.1),(23.1,17.4),(11.8,1.0),(11.4,8.6),(12.8,6.8),(24.1,7.6),(18.8,2.8),(7.0,3.5),(15.4,0.6),(18.7,5.6),(5.2,7.1),(19.2,9.7),(6.8,8.0),(11.9,5.4),(11.5,9.9),(17.8,14.5),(25.5,3.4),(25.9,11.1),(25.8,3.1),(14.0,9.7),(7.0,16.3),(23.8,2.3),(12.0,12.3),(16.9,2.4),(17.1,15.9),(10.3,2.6),(18.5,2.7),(16.5,12.3),(23.0,4.4),(28.0,10.4),(0.9,11.4),(14.1,10.1),(19.2,10.9),(24.6,10.2),(24.0,6.8),(18.8,2.0),(22.3,12.7),(30.1,10.8),(18.8,12.4),(34.9,11.1),(13.5,5.6),(3.9,10.7),(22.4,14.7),(35.8,15.0),(30.1,14.7),(30.5,13.0),(24.9,1.4),(5.7,8.5),(4.2,15.9),(19.0,1.9),(9.8,14.4),(11.3,10.1),(13.2,8.3),(19.5,9.0),(16.2,11.1),(13.2,12.8),(18.5,11.7),(22.7,4.6),(12.9,8.9),(17.3,1.9),(16.8,3.3),(18.4,14.6),(30.6,4.0),(22.8,10.3),(26.0,6.1),(32.1,7.1),(18.0,3.3),(13.0,14.1),(5.5,17.3),(12.2,7.8),(25.8,14.2),(12.7,4.1),(25.8,6.4),(24.5,10.4),(30.7,15.0),(18.0,11.8),(11.2,13.9),(17.9,6.2),(11.7,13.9),(17.9,3.4),(35.8,7.6),(10.2,8.1),(17.3,5.4),(25.3,13.3),(20.0,6.4),(24.5,2.5),(26.1,15.7),(18.4,13.4),(28.7,13.0),(16.2,7.0),(24.4,13.8),(6.3,5.8),(13.6,16.5),(1.0,10.6),(12.3,3.4),(13.5,9.1),(0.5,14.4),(17.7,15.1),(13.6,11.7),(18.5,5.9),(30.2,13.4),(17.6,14.7),(6.9,8.6),(18.5,8.4),(13.6,15.0),(28.7,9.8),(18.2,13.3),(19.0,5.8),(17.4,6.7),(18.3,12.4),(35.3,3.4),(24.4,3.0),(1.2,7.7),(18.6,11.3),(17.6,10.8),(11.3,17.3),(11.9,3.9),(11.3,11.2),(12.3,3.2),(28.0,3.8),(29.6,10.5),(12.1,4.6),(0.8,2.5),(21.5,2.0),(30.1,2.9),(28.4,4.6),(24.4,3.2),(21.7,11.3),(30.0,16.5),(6.1,7.9),(16.6,3.9),(17.0,1.2),(18.7,16.1),(22.7,14.7),(5.4,9.7),(23.7,0.8),(25.3,11.3),(5.2,5.6),(19.7,4.0),(12.7,10.1),(24.8,10.1),(11.5,11.2),(16.8,13.6),(18.5,10.6),(22.4,13.9),(6.0,10.5),(13.5,17.2),(11.6,7.6),(7.4,17.3),(6.7,0.6),(25.2,9.5),(24.1,14.4),(12.4,10.6),(10.5,11.5),(17.2,6.6),(17.9,5.9),(28.8,9.5),(17.9,15.5),(17.0,11.3),(6.9,4.7),(22.3,2.2),(20.6,10.8),(24.1,7.0),(18.4,17.0),(31.4,6.2),(25.1,0.8),(11.8,15.1),(22.9,4.4),(24.3,14.7),(17.4,11.8),(25.4,12.6),(34.4,16.5),(23.6,4.5),(23.6,1.0),(18.3,3.5),(23.7,7.4),(7.7,2.6),(13.0,14.9),(17.9,7.3),(31.4,11.2),(23.5,15.3),(13.6,11.9),(12.1,13.8),(17.8,8.0),(6.2,12.7),(11.7,4.0),(21.9,8.0),(13.1,11.3),(25.4,6.3),(3.9,5.3),(13.1,8.0),(17.5,16.1),(13.9,1.4),(19.4,9.8),(24.6,10.8),(18.6,5.2),(3.8,8.2),(12.8,6.4),(7.6,8.5),(19.5,8.3),(19.2,8.0),(13.0,4.3),(4.4,17.2),(17.1,5.8),(18.5,9.6),(21.8,5.9),(24.5,0.9),(11.1,14.3),(17.4,10.7),(23.0,12.6),(14.5,5.4),(17.2,2.3),(18.3,11.7),(18.3,13.3),(5.8,6.6),(17.9,6.3),(24.8,2.9),(23.6,9.2),(17.3,3.9),(30.1,10.9),(17.4,5.2),(19.5,13.1),(12.2,5.1),(24.0,8.5),(11.0,6.5),(13.3,13.3),(24.7,12.7),(17.1,14.4),(6.7,14.6),(4.6,9.2),(17.9,11.6),(30.6,3.3),(29.5,11.2),(31.0,3.9),(17.4,15.2),(6.8,11.5),(12.3,7.4),(29.7,9.9),(16.6,5.9),(30.5,3.5),(17.1,10.9),(12.1,13.7),(24.0,8.5),(25.0,9.3),(22.7,3.0),(12.2,8.8),(0.8,15.0),(13.2,11.6),(6.9,5.4),(29.1,7.3),(19.4,16.3),(22.0,11.0),(18.3,14.6),(24.5,4.3),(24.2,6.3),(13.7,7.5),(31.7,5.6),(12.8,5.9),(4.0,15.8),(22.5,6.1),(18.3,11.3),(24.2,16.0),(30.8,10.4),(18.3,6.1),(29.2,5.1),(18.4,2.1),(4.7,3.3),(23.7,11.8),(11.0,1.6),(23.9,1.3),(22.9,6.2),(13.4,14.3),(18.5,2.9),(17.5,12.2),(4.8,6.8),(12.7,3.7),(7.1,15.7),(12.8,14.2),(22.9,17.4),(24.7,15.7),(23.6,13.1),(16.6,11.4),(11.4,14.1),(7.7,9.6),(25.9,7.7),(25.0,12.9),(17.3,11.0),(19.2,4.6),(10.9,1.4),(11.9,7.2),(5.7,6.3),(16.7,11.3),(23.1,7.7),(18.6,6.6),(30.4,4.3),(21.7,16.9),(29.3,11.4),(23.1,14.9),(4.6,10.2),(26.2,12.0),(18.0,7.4),(6.0,16.4),(24.9,2.4),(11.2,14.6),(25.7,8.6),(19.5,17.1),(18.1,12.4),(12.4,9.9),(31.3,13.4),(6.0,13.3),(17.8,14.2),(12.7,16.2),(24.1,15.0),(11.7,10.9),(16.2,3.8),(4.7,9.9),(6.1,15.1),(12.7,11.0),(18.1,6.6),(14.0,12.2),(11.0,0.9),(17.7,12.3),(23.0,9.1),(5.1,11.7),(20.4,2.0),(23.1,5.0),(16.4,8.3),(19.0,1.9),(25.3,15.5),(17.7,1.2),(18.4,14.2),(6.8,7.8),(23.5,3.9),(14.2,1.8),(23.9,1.4),(19.1,16.0),(10.5,16.0),(24.3,6.9),(24.3,1.3),(18.3,9.2),(8.0,15.5),(22.6,17.2),(10.6,7.3),(24.8,8.7),(7.6,15.8),(12.9,11.7),(6.6,3.8),(18.3,6.8),(0.3,3.6),(12.5,14.1),(25.5,13.5),(8.1,14.9),(12.6,17.1),(19.2,17.2),(1.8,4.7),(18.4,4.3),(5.7,3.5),(12.8,14.7),(10.3,4.8),(29.1,17.3),(24.5,15.2),(23.8,9.8),(5.5,14.4),(1.8,10.8),(16.9,11.4),(34.3,15.4),(24.8,2.4),(1.1,17.1),(7.6,4.1),(16.3,7.5),(19.5,3.4),(31.5,13.1),(4.5,7.4),(31.4,16.0),(5.9,8.8),(20.1,0.9),(23.8,8.6),(29.8,17.0),(34.2,1.6),(17.1,4.8),(27.4,6.3),(19.0,14.7),(18.8,15.2),(29.7,14.4),(13.0,3.8),(30.1,15.8),(16.4,5.6),(28.8,14.7),(10.4,8.0),(12.8,17.0),(6.7,14.4),(18.3,13.5),(13.9,2.0),(17.7,4.8),(22.5,8.7),(17.8,8.5),(22.7,12.5),(11.9,7.2),(17.9,14.8),(18.3,4.8),(25.8,4.2),(24.2,1.4),(0.5,6.0),(22.9,8.2),(18.5,11.1),(12.1,1.7),(12.2,10.4),(16.2,8.1),(4.7,2.3),(18.4,6.8),(17.7,13.9),(0.4,5.5),(12.3,15.1),(31.2,3.3),(17.0,3.8),(21.7,2.3),(4.6,5.2),(7.0,15.2),(18.6,3.3),(12.1,14.0),(13.3,7.2),(13.5,9.7),(17.9,8.1),(7.2,1.7),(18.6,1.1),(12.9,14.2),(25.2,1.3),(18.3,8.0),(16.9,1.7),(31.1,16.4),(24.6,0.7),(12.4,7.2),(35.5,11.4),(18.4,10.2),(6.2,9.3),(22.2,5.3),(13.5,6.4),(5.7,11.1),(30.6,3.7),(24.5,10.5),(23.0,8.0),(6.0,8.8),(9.9,0.9),(11.5,9.6),(17.0,4.1),(28.6,3.0),(24.8,15.8),(16.4,7.2),(26.4,15.5),(24.3,8.8),(12.2,7.8),(6.8,1.4),(25.6,13.0),(11.6,1.8),(17.5,10.6),(13.1,16.1),(18.9,14.9),(22.3,1.7),(19.0,2.5),(4.0,5.5),(21.9,5.8),(17.7,7.7),(13.4,1.9),(28.6,10.7),(17.5,12.5),(6.8,1.8),(31.5,5.7),(5.8,2.6),(11.7,12.6),(13.9,9.3),(13.2,17.1),(22.9,3.0),(12.0,11.4),(34.4,3.1),(19.5,2.0),(10.0,11.7),(29.5,4.4),(15.9,8.3),(12.4,15.2),(30.5,4.0),(23.5,9.4),(10.7,1.0),(30.8,16.8),(2.0,6.2),(17.7,3.2),(23.6,5.0),(5.3,5.1),(2.0,15.7),(11.6,12.5),(29.9,6.7),(5.8,10.6),(25.2,11.9),(28.9,13.6),(18.2,3.5),(19.2,0.9),(14.3,2.2),(20.0,1.1),(12.9,8.1),(7.5,15.3),(12.8,15.1),(17.7,5.3),(35.8,17.0),(35.7,6.6),(17.5,10.1),(17.9,5.6),(17.1,8.2),(6.2,8.6),(30.6,12.1),(8.8,17.1),(31.3,10.4),(18.0,5.3),(23.1,5.9),(30.2,10.8),(19.4,14.4),(6.8,9.4),(17.1,5.6),(30.7,16.6),(29.9,5.3),(24.5,9.5),(16.5,7.8),(12.1,0.9),(18.4,15.9),(12.4,11.6),(4.7,13.6),(17.4,16.3),(18.8,13.2),(16.4,7.6),(25.3,8.8),(17.7,9.3),(2.0,9.8),(18.8,3.7),(28.3,3.6),(18.9,9.7),(23.4,5.4),(1.3,16.8),(12.9,11.8),(6.7,16.0),(18.7,5.5),(29.3,10.0),(19.1,5.7),(13.2,8.6),(24.3,11.1),(15.9,3.9),(29.2,12.5),(14.1,9.3),(8.1,10.5),(11.9,15.2),(28.0,15.9),(22.7,7.6),(10.4,9.3),(25.7,2.8),(19.6,10.7),(13.2,14.0),(12.9,14.7),(18.3,4.2),(24.9,1.6),(22.3,12.0),(19.9,7.0),(23.9,10.9),(12.4,8.8),(17.7,10.1),(16.7,13.0),(28.6,12.5),(16.9,1.0),(24.3,14.5),(17.5,9.9),(10.1,6.1),(23.6,11.7),(7.1,14.2),(5.7,14.3),(25.8,16.5),(17.1,7.2),(29.7,11.3),(16.1,5.9),(5.6,4.2),(18.3,6.0),(25.2,0.7),(13.3,6.5),(11.5,13.2),(31.3,8.8),(6.8,1.7),(30.2,12.7),(24.6,15.4),(12.0,15.8),(13.4,14.4),(28.8,4.4),(13.0,8.0),(22.6,14.7),(0.5,3.1),(11.5,1.7),(16.2,14.4),(22.8,10.7),(10.4,0.8),(16.7,4.5),(17.5,13.5),(6.1,9.3),(25.2,0.7),(34.6,5.3),(24.0,15.7),(14.5,4.3),(24.8,5.5),(5.4,7.7),(19.8,6.2),(11.0,7.5),(12.4,17.4),(12.4,12.1),(7.8,15.5),(16.5,0.8),(12.5,2.5),(19.2,6.2),(23.4,4.2),(1.3,4.4),(23.0,8.5),(18.1,13.9),(29.5,3.1),(16.5,9.9),(1.0,10.4),(18.4,2.6),(13.3,3.3),(15.8,5.6),(12.7,16.8),(29.2,1.2),(17.3,1.0),(23.6,14.4),(28.3,7.8),(19.5,8.6),(17.7,2.8),(22.7,10.2),(31.3,7.1),(12.9,10.9),(29.2,16.0),(10.7,4.9),(16.4,4.3),(18.5,5.2),(23.4,11.8),(30.6,9.8),(29.5,15.2),(24.2,3.1),(16.7,13.7),(18.3,11.2),(22.8,4.8),(12.9,17.1),(28.2,7.5),(13.7,15.7),(30.4,7.8),(24.2,4.4),(16.4,9.7),(19.9,2.5),(18.0,7.3),(18.6,1.2),(12.5,5.2),(10.3,6.1),(13.7,12.0),(11.7,12.4),(25.7,14.0),(0.9,0.8),(12.3,15.5),(6.6,6.8),(6.8,12.4),(5.2,9.8),(23.5,4.1),(4.6,6.8),(11.7,13.5),(10.4,9.5),(23.9,17.4),(18.3,6.4),(24.7,6.4),(23.1,4.8),(7.9,14.0),(18.1,7.1),(16.4,6.4),(4.8,7.5),(1.1,0.8),(25.6,3.0),(23.3,11.5),(18.2,5.3),(25.3,3.9),(23.1,5.3),(10.4,10.8),(16.1,15.1),(11.7,5.9),(6.1,13.9),(10.9,4.3),(11.9,7.9),(16.4,16.5),(11.0,14.3),(16.6,6.5),(22.0,16.1),(1.5,2.8),(12.6,8.4),(23.4,15.1),(23.2,7.3),(23.8,14.2),(35.9,1.5),(28.8,12.8),(23.0,12.8),(22.8,12.3),(28.3,6.0),(7.4,13.4),(12.5,14.6),(17.4,14.5),(5.7,15.0),(11.5,14.5),(18.2,13.4),(24.9,16.2),(18.4,4.0),(23.1,15.9),(18.5,13.5),(23.2,14.7),(11.9,4.0),(22.9,12.9),(18.3,13.1),(18.6,13.8),(12.8,16.8),(29.4,4.1),(24.1,16.9),(24.2,8.4),(19.4,10.1),(12.7,6.4),(23.4,7.3),(22.7,12.4),(13.1,3.2),(19.5,3.9),(26.0,8.7),(19.0,11.5),(23.3,6.0),(19.3,2.9),(18.5,3.4),(6.6,13.1),(18.5,4.9),(12.2,12.0),(24.5,6.7),(19.6,13.0),(12.0,10.6),(18.2,11.3),(15.6,3.3),(5.7,14.5),(17.3,5.2),(14.1,17.0),(15.9,15.8),(12.5,12.3),(25.1,1.0),(12.3,15.5),(22.8,8.9),(23.0,1.9),(18.8,7.9),(11.3,9.1),(19.2,11.5),(24.1,1.4),(18.7,8.5),(7.2,6.8),(17.6,10.4),(22.3,5.4),(10.4,9.3),(23.2,3.5),(24.4,14.6),(12.8,2.8),(11.2,15.4),(11.8,4.9),(35.0,14.1),(19.1,8.7),(12.8,12.4),(23.1,7.8),(12.6,16.5),(17.4,11.1),(1.6,6.9),(11.9,16.9),(24.0,2.8),(17.8,13.7),(14.3,16.0),(11.1,15.4),(13.8,15.9),(25.2,11.1),(28.1,3.8),(19.8,5.5),(23.3,8.9),(18.4,6.8),(18.3,7.7),(28.8,5.3),(16.4,12.7),(0.5,8.1),(6.1,8.1),(23.8,8.1),(18.5,2.6),(24.1,5.6),(24.0,16.2),(10.8,13.5),(10.6,4.7),(23.6,0.8),(16.5,10.2),(6.6,5.7),(24.6,17.0),(25.2,8.1),(17.8,15.6),(11.4,1.3),(30.0,6.5),(23.2,9.4),(19.0,7.5),(23.8,13.2)} {\fill[bkBlue!80!black] \p circle (0.25);}
  \node[txc lab, anchor=north] at (18,-1) {光子 1000 個};
  \end{scope}
\end{tikzpicture}
\bkfigcaption{光子を1個ずつ数えたときの、二重スリットのスクリーン（Python の乱数で、干渉じまの強さの分布にしたがって点を打った模式図）。1個1個はばらばらの所に届くが、数が増えると、しまが現れる。}\label{ph:原-2:fig-ichiko}
\end{bkfigure}

\noindent 光子1個がどこに届くかは、前もって決まっていません。決まっているのは、「どこに届きやすいか」という\textbf{確率}で、その確率の分布が、波の干渉の強さの分布と同じなのです。光は「波か粒か」のどちらかではなく、届くときは粒として1個ずつ、どこに届きやすいかは波として、ふるまいます。同じことが、電子でも起こります（原-3）。

\begin{phmijika}[人の目は、光子を何個から感じる？]
暗い所に30分ほどいて目を慣らすと、目のおくの、暗さに強い細胞（かん体）は、光子1個にも反応します。1942年のヘクトたちの実験では、とても弱い光を目に当て、数個〜十数個の光子が細胞に届くと「光った」と感じることがわかりました。わたしたちの目は、光のつぶを数えられるほど感度のよいセンサーなのです。
\end{phmijika}

\phkokomade

\Needspace{10\baselineskip}\phstage{shiki}{光のモードを数えて導く}{5}

\bksubsection{9}{箱の中の光のモードを数える}
\label{ph:原-2:sec9}
ここからは、はてな？の「熱した物の光の分布」を、光子の考えから導きます。黒体の光を調べるには、内側の壁が同じ温度 $T$ に保たれた箱（空どう）を考え、小さな穴から出てくる光を見るのがいちばんです。穴に入った光は、中で吸われてしまうので、穴は「光を全部吸う黒体」になり、穴から出る光は黒体の光です（熱-13 の黒体炉）。

箱の中の光は、壁ではね返って、定在波になります。『波動』（波-6）の弦と同じで、1辺 $L$ の箱の中にできる定在波は、壁の間に半波長がちょうど整数個入るものだけです。$x$ の向きに進む波では、$L=n_x\lambda/2$ なので、波数 $k=2\pi/\lambda$ は
\[ k_x=\frac{\pi}{L}n_x\qquad(n_x=1,\ 2,\ 3,\ \ldots) \]
に限られます。3次元の箱では、3つの向きの整数の組 $(n_x,\ n_y,\ n_z)$ ごとに1つの定在波（\textbf{モード}）ができ、その波数の大きさは
\[ k=\frac{\pi}{L}\sqrt{n_x^2+n_y^2+n_z^2} \]
です。光には、ゆれの向きが2通り（偏り：波-13）あるので、1つの組に2つのモードがあります。

\begin{bkfigure}
\begin{tikzpicture}[x=1mm, y=1mm]
  \begin{scope}
  \draw[bkInk, line width=1pt] (0,-2) -- (0,34) (40,-2) -- (40,34);
  \foreach \nn/\yy in {1/28,2/16,3/4} {
    \draw[bkBlue!80!black, line width=1pt, domain=0:40, samples=60] plot (\x,{\yy+4*sin(\x/40*180*\nn)});
    \draw[bkBlue!40, line width=0.6pt, dash pattern=on 1.5pt off 1.5pt, domain=0:40, samples=60] plot (\x,{\yy-4*sin(\x/40*180*\nn)});
    \node[txc lab, anchor=west] at (41.5,\yy) {$n=\nn$};}
  \draw[txc dim] (0,-5) -- (40,-5); \node[txc lab, anchor=north] at (20,-6) {箱の1辺 $L$};
  \end{scope}
  \begin{scope}[shift={(66,-4)}]
  \foreach \i in {1,...,7} {\foreach \j in {1,...,7} {\fill[bkInk] ({\i*5},{\j*5}) circle (0.55);}}
  \draw[bkInk, line width=0.6pt, -{Stealth[length=2mm]}] (0,0) -- (40,0); \draw[bkInk, line width=0.6pt, -{Stealth[length=2mm]}] (0,0) -- (0,40);
  \node[txc lab, anchor=north] at (38,-0.5) {$n_x$}; \node[txc lab, anchor=east] at (-0.5,38) {$n_y$};
  \draw[bkOrange!85!black, line width=1.2pt] (25,0) arc (0:90:25);
  \draw[bkOrange!85!black, line width=0.8pt, dash pattern=on 2pt off 1.5pt] (28,0) arc (0:90:28);
  \node[txc lab, anchor=west, text width=26mm, align=left] at (42,22) {半径の中の点の数が、その振動数より低いモードの数};
  \end{scope}
\end{tikzpicture}
\bkfigcaption{左：1辺 $L$ の箱にできる定在波（弦の固有振動と同じ。$n$ は半波長の数）。右：整数の組 $(n_x,\ n_y)$ を点で表す（2次元の場合）。決まった振動数より低いモードは、原点を中心とする円（3次元では球）の中の点。}\label{ph:原-2:fig-mode}
\end{bkfigure}

\phboxhead{球の体積で数える}
振動数は $f=ck/(2\pi)$（波-5）なので、振動数 $f$ より低いモードは、$(n_x,\ n_y,\ n_z)$ の点のうち、半径 $\sqrt{n_x^2+n_y^2+n_z^2}<2Lf/c$ の球の中にあるものです。$n$ はどれも正の数なので、球の8分の1だけを数えます。点は、1辺1の立方体に1個ずつ並んでいるので、点の数は体積で数えられます（4-5 の球の体積 $\tfrac43\pi r^3$）。偏りの2通りをかけると、
\[ N=2\times\frac18\times\frac43\pi\left(\frac{2Lf}{c}\right)^3=\frac{8\pi f^3}{3c^3}L^3 \]
です。箱の体積 $L^3$ でわり、$f$ で微分すると、\qty{1}{m^3} あたり、振動数 $f$ から $f+\dd f$ までのはばに入るモードの数は
\pheq{\frac{8\pi f^2}{c^3}\,\dd f}{1/m$^3$ ＝ Hz$^2$/(m/s)$^3$ × Hz}
\label{ph:原-2:modes}です。振動数が高いほど、モードは $f^2$ に比例して、どんどん多くなります。たとえば、\qty{1}{cm^3} の箱の中で、波長 \qty{550}{nm} のまわりの \qty{1}{nm} のはばに入るモードは、約 $2.7\times10^{11}$ 個もあります。

\bksubsection{10}{モード1つの平均のエネルギー：紫外の破綻}
\label{ph:原-2:sec10}
次に、1つのモードが、温度 $T$ のとき平均でどれだけのエネルギーをもつかを考えます。

\phboxhead{波だけで考えると：どのモードも $k_{\mathrm B}T$}
光の1つのモードは、ばねのゆれと同じ形の振動です。熱-14・熱-15 の\textbf{エネルギー等分配の法則}では、ゆれ1つは平均で $k_{\mathrm B}T$ のエネルギーをもちます。これをそのまま使うと、\qty{1}{m^3}・振動数のはば $\dd f$ あたりの光のエネルギー $u_f\,\dd f$ は
\[ u_f=\frac{8\pi f^2}{c^3}k_{\mathrm B}T \]
になります（レイリー・ジーンズの式）。この式は、振動数の低い所（赤外線より長い波長）では測定と合います。ところが、振動数が高くなるほど $f^2$ でどこまでも大きくなり、全部の振動数でたすと、エネルギーが無限大になってしまいます。熱した物は、紫外線や X 線を、とてつもなく強く出すことになるのです。実際にはそうならないこのくいちがいを、\textbf{紫外の破綻}（紫外発散）といいます。\bkidx{しがいのはたん}{紫外の破綻}

\phboxhead{エネルギーが $hf$ ずつなら}
プランクの仮定では、振動数 $f$ のモードのエネルギーは、$hf$ の何個分か（$0,\ hf,\ 2hf,\ \ldots$）にしかなれません。これは、熱-15 ⑧の「エネルギーの粒 $\Delta\varepsilon$」の考えで $\Delta\varepsilon=hf$ としたものと同じです。そこでは、ボルツマン因子の和が等比級数（7-10⑤）になり、平均のエネルギーは
\pheq{\langle\varepsilon\rangle=\frac{hf}{e^{hf/(k_{\mathrm B}T)}-1}}{J ＝ J}
\label{ph:原-2:avg}でした。光子の数で言えば、モード1つに平均 $1/(e^{hf/(k_{\mathrm B}T)}-1)$ 個の光子がいます。これは、熱-15 ⑨のボース分布（化学ポテンシャル0）そのものです。

\begin{bkfigure}
\begin{tikzpicture}
\begin{axis}[bkplot, width=96mm, height=48mm, xmin=0, xmax=8.2, ymin=0, ymax=1.15, xtick={1,2,3,4,5,6,7,8}, ytick={0.25,0.5,0.75,1}, yticklabel style={/pgf/number format/fixed}, xlabel={$x=hf/(k_{\mathrm B}T)$}, ylabel={$\langle\varepsilon\rangle/(k_{\mathrm B}T)$}]
  \addplot[bkthin, dashed, line width=1.2pt, domain=0:8] {1};
  \addplot[bkcurve, domain=0.01:8, samples=120] {x/(exp(x)-1)};
  \fill[bkRed] (axis cs:4.53,0.0493) circle (1.8pt);
  \node[font=\sffamily\footnotesize, anchor=south west, text=bkRed!75!black] at (axis cs:4.6,0.07) {太陽の \qty{550}{nm}（例題11）};
  \node[font=\sffamily\footnotesize, anchor=south east] at (axis cs:8,1.02) {破線：波だけで考えたとき（いつも $k_{\mathrm B}T$）};
  \node[font=\sffamily\footnotesize, anchor=west, text=bkBlue!80!black] at (axis cs:1.2,0.55) {$hf$ ずつのとき};
\end{axis}
\end{tikzpicture}
\bkfigcaption{モード1つの平均のエネルギー。光子1個のエネルギー $hf$ が熱のエネルギー $k_{\mathrm B}T$ より大きいモードは、ほとんどエネルギーをもてない。}\label{ph:原-2:fig-avg}
\end{bkfigure}

\begin{phyomi}[$\langle\varepsilon\rangle=\dfrac{hf}{e^{hf/(k_{\mathrm B}T)}-1}$]
「1つぶが大きすぎるモードには、熱のエネルギーが入りにくい」。$hf\ll k_{\mathrm B}T$（振動数の低いモード）では $e^x\approx1+x$（5-17）から $\langle\varepsilon\rangle\approx k_{\mathrm B}T$ で、波だけの考えと同じです。$hf\gg k_{\mathrm B}T$ では $\langle\varepsilon\rangle\approx hf\,e^{-hf/(k_{\mathrm B}T)}$ で、急に0に近づきます。モードの数は $f^2$ で増えても、1つあたりのエネルギーが指数関数で減るので、高い振動数の光は出ないのです。\textbf{これが、はてな？の答えのかぎです}。
\end{phyomi}

\bksubsection{11}{プランクの式・ウィーンの変位則・シュテファン・ボルツマン定数}
\label{ph:原-2:sec11}
モードの数（⑨）に、モード1つの平均のエネルギー（⑩）をかけると、箱の中の \qty{1}{m^3}・はば $\dd f$ あたりの光のエネルギーは
\pheq{u_f\,\dd f=\frac{8\pi hf^3}{c^3}\,\frac{1}{e^{hf/(k_{\mathrm B}T)}-1}\,\dd f}{J/m$^3$}
\label{ph:原-2:planck}です。これが\textbf{プランクの式}です。\bkidx{ぷらんくのしき}{プランクの式}熱-13 のように波長で書くには、$f=c/\lambda$、$|\dd f|=c\,\dd\lambda/\lambda^2$ で書きかえます。さらに、穴から出ていく光は、箱の中の光のエネルギーの $c/4$ 倍（いろいろな向きに進む光を平均した値。熱の専門の本で導く）なので、黒体の表面 \qty{1}{m^2} から出る、波長 $\lambda$ のまわりの単色放射の強さは
\[ E_{\mathrm b\lambda}=\frac{2\pi hc^2}{\lambda^5}\,\frac{1}{e^{hc/(\lambda k_{\mathrm B}T)}-1} \]
です。熱-13 の式とくらべると、測定から決めた2つの定数が、プランク定数で書けたことになります。
\[ c_1=2\pi hc^2=\qty{3.742e-16}{W.m^2},\qquad c_2=\frac{hc}{k_{\mathrm B}}=\qty{1.4388e-2}{m.K} \]
$c_1$ を W·\unit{\micro m}$^4$/m$^2$ で書くと $3.742\times10^8$、$c_2$ を \unit{\micro m}·K で書くと \num{14388} で、熱-13 の値と同じです（例題12）。

\phboxhead{山の位置：ウィーンの変位則}
$x=hc/(\lambda k_{\mathrm B}T)$ とおくと、$E_{\mathrm b\lambda}$ は $x^5/(e^x-1)$ に比例します。山の位置は、微分して0になる所（5-7）です。
\[ \frac{\dd}{\dd x}\,\frac{x^5}{e^x-1}=\frac{5x^4(e^x-1)-x^5e^x}{(e^x-1)^2}=0\qquad\text{より}\qquad x=5(1-e^{-x}) \]
この式は、$x=5$ から始めて、右辺に入れて左辺を出す、をくり返すと、4.966、4.965、4.965 とすぐに落ち着きます。$x=4.965$ なので、
\pheq{\lambda_{\max}T=\frac{hc}{4.965\,k_{\mathrm B}}=\qty{2.898e-3}{m.K}}{m·K}
\label{ph:原-2:wien}です。熱-13 のウィーンの変位則が、$h$・$c$・$k_{\mathrm B}$ から出てきました。

\phboxhead{全部の波長でたす：シュテファン・ボルツマン定数}
$E_{\mathrm b\lambda}$ を全部の波長で積分すると、$x$ に変数を変えて
\[ E_{\mathrm b}=\int_0^\infty E_{\mathrm b\lambda}\,\dd\lambda=\frac{2\pi k_{\mathrm B}^4T^4}{h^3c^2}\int_0^\infty\frac{x^3}{e^x-1}\,\dd x \]
となります。積分の値は、細い長方形に分けてたす（5-10）と 6.494 で、これは $\pi^4/15$ と同じ数です（$\pi^4/15$ になることは、専門の本で無限級数を使って示す）。したがって
\pheq{E_{\mathrm b}=\sigma T^4,\qquad \sigma=\frac{2\pi^5k_{\mathrm B}^4}{15h^3c^2}=\qty{5.67e-8}{W/(m^2.K^4)}}{W/m$^2$}
\label{ph:原-2:sigma}です。熱-13 で「測定で決まる定数」としたシュテファン・ボルツマン定数も、$h$・$c$・$k_{\mathrm B}$ で書けました。

\Needspace{16\baselineskip}\begin{phflow}[プランクの式から、熱の放射の式へ]
\phflowrow{$\langle\varepsilon\rangle=\dfrac{hf}{e^{hf/(k_{\mathrm B}T)}-1}$}{振動数が低い\\$hf\ll k_{\mathrm B}T$}{$\langle\varepsilon\rangle=k_{\mathrm B}T$}
\phflowrow{$E_{\mathrm b\lambda}=\dfrac{c_1}{\lambda^5(e^{c_2/(\lambda T)}-1)}$}{山の位置\\（微分して0）}{$\lambda_{\max}T=\qty{2.898e-3}{m.K}$}
\phflowrow{$E_{\mathrm b\lambda}=\dfrac{c_1}{\lambda^5(e^{c_2/(\lambda T)}-1)}$}{全部の波長で\\積分する}{$E_{\mathrm b}=\sigma T^4$}
\end{phflow}

\begin{phbaisuru}{温度 $T$}
山の波長は半分に、全体の明るさは $2^4=16$ 倍になります。図\ref{ph:原-2:fig-bunpu}の山が左へ動き、背が高くなるのは、このためです。
\end{phbaisuru}

\Needspace{6\baselineskip}\phboxhead{\bkhatenaref の答え}
\label{ph:原-2:hatena}熱した物の光は、光子 $hf$ のつぶで出ていきます。温度 $T$ の物がもつ熱のエネルギーの目安は $k_{\mathrm B}T$ で、$hf$ が $k_{\mathrm B}T$ の数倍より大きい光子は、ほとんど出せません（⑩）。\qty{1000}{K} では $k_{\mathrm B}T=\qty{0.086}{eV}$ で、赤い光の光子（約 \qty{1.9}{eV}）でも約22倍、青い光の光子（約 \qty{2.8}{eV}）は約32倍です。どちらもめったに出せませんが、赤のほうがずっと出やすいので、暗い赤に光ります（全体の光の約 0.001\,\% が可視光）。温度が上がると、大きなつぶの光子も出せるようになり、黄色・緑・青の光が加わって、白に近づくのです（\qty{1500}{K} で約 0.17\,\%、\qty{2800}{K} で約10\,\%、\qty{5772}{K} で約47\,\%：例題1）。

\begin{motto}[4]
\phboxhead{宇宙は、温度 \qty{2.7}{K} の黒体の光で満ちている}
宇宙のどの方向からも、とても弱い電波（マイクロ波）がやって来ます。宇宙が生まれてまもない、熱くて明るかったころの光が、宇宙が広がるにつれて波長がのびたものだと考えられていて、\textbf{宇宙マイクロ波背景放射}とよばれます。1990年ごろ、人工衛星でこの光の分布を正確にはかると、温度 約 \qty{2.725}{K} の黒体のプランクの式（⑪）に、とてもよく合いました（2006年のノーベル物理学賞）。山の波長は、ウィーンの変位則から $\qty{2.898e-3}{m.K}\div\qty{2.725}{K}\approx\qty{1.06}{mm}$ です。プランクの式は、実験室の箱の中だけでなく、宇宙全体の光にも当てはまるのです。
\end{motto}

\bksubsection{12}{光センサー・放射温度計・暗視カメラ}
\label{ph:原-2:sec12}
光子の考えは、光をはかる機械の設計の土台です。
\begin{itemize}[leftmargin=1.2\zw, itemsep=2pt, topsep=2pt, label={\color{bkTheme}●}]
\item \textbf{光電子増倍管}：光電面から飛び出した光電子を、電圧で加速して次々に別の電極に当て、電子の数を100万倍ほどに増やしてはかる。光子1個でも数えられ、ニュートリノの観測装置スーパーカミオカンデ（2015年のノーベル物理学賞の研究に使われた）の壁には、直径 \qty{50}{cm} のものが約1万本並んでいる。
\item \textbf{フォトダイオードとイメージセンサー}：半導体の中で、光子1個が電子1個を動けるようにする（内部の光電効果）。光子のエネルギーが半導体の禁制帯の幅より大きくないと感じない（例題6・原-9）。スマホのカメラの1つ1つの画素は、とどいた光子の数を数えている。
\item \textbf{放射温度計・暗視カメラ}：物が出す赤外線の強さと分布から、さわらずに温度をはかる（熱-13）。人の体（約 \qty{310}{K}）の光の山は約 \qty{9.3}{\micro m} で、暗視カメラ（サーモグラフィー）は、大気をよく通る \qtyrange{8}{14}{\micro m} の赤外線（人の体の光の約38\,\%）を見ている。
\end{itemize}

\begin{phmada}[光が当たってから、電子が飛び出すまでの時間は？]
光電効果では、光が当たると「すぐに」電子が飛び出します（②）。では、その「すぐ」は、どれくらい短いのでしょう。2000年代になって、アト秒（$10^{-18}$ 秒）の光のパルスを作れるようになり（2023年のノーベル物理学賞）、電子が飛び出す時間の、わずかな差がはかれるようになりました。2010年の実験では、ネオンの原子の中の、異なる場所（軌道：原-7）にいる電子が、同じ光で飛び出すとき、約20アト秒の時間の差があることがわかりました。この差は、飛び出す電子と、原子に残った電子どうしがおよぼし合う力が関係すると考えられていますが、計算と測定がどこまで合うのかは、研究が続いています。
\end{phmada}

\begin{phsaki}{量子光学と光をはかる技術}
この章では、光子を「つぶ」として数えました。光子を1個ずつ出したり、2つの光子を結びつけたりする技術（量子光学）は、盗み見のできない暗号の通信や、量子コンピュータの研究に使われています（2022年のノーベル物理学賞は、光子どうしの不思議なつながり「量子もつれ」の実験）。光検出器の設計、放射による温度の測定（放射測温）は、光エレクトロニクスや計測工学の専門の本で学びます。
\end{phsaki}

\begin{senmonsho}[物理の専門書の書き方]
\begin{bkhonbun}
光子のエネルギーは $E=h\nu=\hbar\omega$，運動量は $\symbfit{p}=\hbar\symbfit{k}$ である．空どう放射のエネルギー密度は
\[ u(\nu)\,\dd\nu=\frac{8\pi h\nu^3}{c^3}\frac{\dd\nu}{e^{h\nu/k_{\mathrm B}T}-1} \]
\end{bkhonbun}
\bkyaku{$\nu$ はこの本の $f$（振動数）。$\hbar=h/(2\pi)$（エイチバー）、$\omega=2\pi f$ は角振動数なので $\hbar\omega=hf$。$\symbfit{k}$ は波数ベクトル（大きさ $2\pi/\lambda$）で、$\hbar k=h/\lambda$。$u(\nu)$ は⑪の $u_f$ と同じ。}
\end{senmonsho}

\begin{taisetsu}
\def\tsl#1{\makebox[7.6\zw][l]{\sffamily\bfseries\small\color{\bkTT}#1}}
\tsl{光子}$E=hf=\dfrac{hc}{\lambda}$、$hc=\qty{1240}{eV.nm}$、$p=\dfrac{h}{\lambda}=\dfrac{E}{c}$\\[3pt]
\tsl{光電効果}$K_0=hf-W$、$f_0=\dfrac{W}{h}$、$eV_0=K_0$\\[3pt]
\tsl{コンプトン}$\Delta\lambda=\dfrac{h}{m_{\mathrm e}c}(1-\cos\theta)$、$\dfrac{h}{m_{\mathrm e}c}=\qty{2.43}{pm}$\\[3pt]
\tsl{モードの数}\qty{1}{m^3}・$\dd f$ あたり $\dfrac{8\pi f^2}{c^3}\dd f$\\[3pt]
\tsl{モードの平均}$\langle\varepsilon\rangle=\dfrac{hf}{e^{hf/(k_{\mathrm B}T)}-1}$\\[3pt]
\tsl{熱の放射}$c_1=2\pi hc^2$、$c_2=\dfrac{hc}{k_{\mathrm B}}$、$\lambda_{\max}T=\dfrac{hc}{4.965k_{\mathrm B}}$、$\sigma=\dfrac{2\pi^5k_{\mathrm B}^4}{15h^3c^2}$
\end{taisetsu}

\begin{kotoba}
\small
\bkword{光子}{エネルギー $hf$・運動量 $h/\lambda$ をもつ光のつぶ。光量子ともいう。}
\bkword{プランク定数}{光子のエネルギーと振動数の比 $h$。$\qty{6.63e-34}{J.s}$。}
\bkword{光電効果}{光を当てた金属から電子が飛び出すこと。}
\bkword{仕事関数}{金属の中の電子が外へ出るのに要る、いちばん小さいエネルギー $W$。}
\bkword{限界振動数}{光電効果が起こる、いちばん低い振動数 $f_0=W/h$。}
\bkword{阻止電圧}{いちばん速い光電子も止める逆向きの電圧 $V_0$。}
\bkword{コンプトン効果}{X 線やγ線が電子にぶつかり、波長がのびること。}
\bkword{黒体・モード}{光を全部吸う理想の物。箱の中にできる1つ1つの定在波。}
\bkword{プランクの式}{黒体の光の強さの分布。モードの数×平均のエネルギー。}
\bkword{紫外の破綻}{波だけの考えでは、高い振動数の光のエネルギーが無限大になること。}
\end{kotoba}

\bkkigocards{\bkkigocard{$h$}{エイチ}{プランク定数〔J·s〕。光子 $E=hf$（教科書は $h\nu$）。高さ・比エンタルピーの $h$ とは別}{Planck constant}{原-2}}
{\bkkigocard{$W$}{ダブリュー}{仕事関数〔eV〕。金属から電子を取り出すのに要るエネルギー。仕事 $W$（第1・2巻）とは別}{work function}{原-2}}
\bkkigocards{\bkkigocard{$K_0$}{ケー・ゼロ}{光電子の運動エネルギーの最大値〔J・eV〕。$K_0=hf-W$}{maximum kinetic energy}{原-2}}
{\bkkigocard{$V_0$}{ブイ・ゼロ}{阻止電圧〔V〕。$eV_0=K_0$。V で書いた数は $K_0$ を eV で書いた数と同じ}{stopping voltage}{原-2}}

%% ======================== 例題 ========================
\Needspace{10\baselineskip}\bksection{reidai}{例題でたしかめよう}
\label{ph:原-2:ex1}
\begin{phreidai}[熱した物の光のうち、目に見えるのはどれだけ？]
\qty{1000}{K}・\qty{1500}{K}・\qty{2800}{K}・\qty{5772}{K}の黒体について、光の山の波長と、出る光全体のうち可視光（\qtyrange{380}{780}{nm}）が占める割合を求めましょう。
\phbamen
はてな？の「熱した鉄の色」を、プランクの式で数にする場面です。
\phmushi
熱した物を黒体とみます（本物の鉄は放射率が1より小さく、波長でも少し変わる：熱-13）。
\phzu
図\ref{ph:原-2:fig-bunpu}。黄色の帯の面積÷曲線の下の全体の面積。
\phriyu
山の波長はウィーンの変位則（⑪）。可視光の部分の面積は、プランクの式 $E_{\mathrm b\lambda}$ を \qtyrange{380}{780}{nm} で積分する。式では積分できないので、細い台形に分けてたす（5-10）。全体の面積は $\sigma T^4$（⑪）。
\phtateru
\[ \lambda_{\max}=\frac{\qty{2.898e-3}{m.K}}{T},\qquad(\text{割合})=\frac{1}{\sigma T^4}\int_{380\,\mathrm{nm}}^{780\,\mathrm{nm}}E_{\mathrm b\lambda}\,\dd\lambda \]
\phtoku
下のプログラムで、\qty{10}{nm}ずつの台形40本でたした。
\begin{phcode}[Python]
from math import exp, pi
h, c, k = 6.626e-34, 2.998e8, 1.381e-23
sigma = 5.670e-8
def E(lam, T):     \# 単色放射の強さ
    x = h * c / (lam * k * T)
    return 2 * pi * h * c**2 / lam**5 / (exp(x) - 1)
for T in [1000, 1500, 2800, 5772]:
    vis = 0.0
    for i in range(40):     \# 380〜780 nm を 10 nm ずつ
        a = 380e-9 + i * 10e-9
        vis = vis + (E(a, T) + E(a + 10e-9, T)) / 2 * 10e-9
    peak = round(2.898e-3 / T * 1e6, 2)   \# 山の波長（μm）
    ratio = vis / (sigma * T**4) * 100    \# 可視光の割合（\%）
    print(T, peak, '\%.2g' \% ratio)
\end{phcode}
\begin{phcodeout}[実行結果]
1000 2.9 0.0011
1500 1.93 0.17
2800 1.04 9.8
5772 0.5 47
\end{phcodeout}
2列めが山の波長（\unit{\micro m}）、3列めが可視光の割合（\%）。
\phtashikame
\qty{2800}{K}の約10\,\%、\qty{5772}{K}の約47\,\% は、熱-13 のもっと知りたいの値と同じ。台形を \qty{1}{nm} ずつにしても、2けたまでは変わらない。\\
大きさ：\qty{1000}{K}では10万分の1ほどしか目に見える光がないので、暗い赤にしか見えない。
\kotae 山の波長 \qty{2.9}{\micro m}・\qty{1.9}{\micro m}・\qty{1.0}{\micro m}・\qty{0.50}{\micro m}。可視光の割合 約 0.0011\,\%・0.17\,\%・10\,\%・47\,\%
\end{phreidai}
\begin{jibun}
例題1のプログラムの温度を変えて、\qty{1300}{K}と\qty{2000}{K}の黒体の、山の波長と可視光の割合を求めましょう。
\bkjans{\qty{1300}{K}：山 \qty{2.2}{\micro m}、約 0.038\,\%。\qty{2000}{K}：山 \qty{1.4}{\micro m}（1.45）、約 1.7\,\%}
\end{jibun}

\label{ph:原-2:ex2}
\begin{phreidai}[光の色と、光子1個のエネルギー]
波長 \qty{650}{nm}（赤）・\qty{550}{nm}（緑）・\qty{450}{nm}（青）・\qty{300}{nm}（紫外線）の光子1個のエネルギーを、eV と J で求めましょう。
\phbamen
色の違う光の、つぶの大きさをくらべる場面です。
\phmushi
波長は空気の中でも真空の中でも同じとします（違いは0.03\,\% ほど）。
\phzu
波長 $\lambda$、光子のエネルギー $E$。
\phriyu
$E=hc/\lambda$（③）。$hc=\qty{1240}{eV.nm}$ を使うと、nm の数でわるだけで eV が出る。
\phtateru
\[ E\ [\mathrm{eV}]=\frac{1240}{\lambda\ [\mathrm{nm}]},\qquad E\ [\mathrm J]=E\ [\mathrm{eV}]\times1.60\times10^{-19} \]
\phtoku
\[ \frac{1240}{650}\approx\qty{1.91}{eV},\quad\frac{1240}{550}\approx\qty{2.25}{eV},\quad\frac{1240}{450}\approx\qty{2.76}{eV},\quad\frac{1240}{300}\approx\qty{4.13}{eV} \]
J では、$hc/\lambda$ に直接入れて、$\qty{3.06e-19}{J}$・$\qty{3.62e-19}{J}$・$\qty{4.42e-19}{J}$・$\qty{6.63e-19}{J}$。
\phtashikame
eV の数に $1.60\times10^{-19}$ をかけても、J の値は3けためが1ずれる程度で合う（丸めのため）。\\
大きさ：可視光は \qtyrange{1.6}{3.3}{eV}。紫外線の光子は可視光より大きく、日焼けを起こす（細胞の分子をこわせる）。
\kotae \qty{1.91}{eV}・\qty{2.25}{eV}・\qty{2.76}{eV}・\qty{4.13}{eV}（$\qty{3.06e-19}{J}$・$\qty{3.62e-19}{J}$・$\qty{4.42e-19}{J}$・$\qty{6.63e-19}{J}$）
\end{phreidai}
\begin{jibun}
ブルーレイの読みとりに使う \qty{405}{nm} の光と、テレビのリモコンの赤外線 \qty{940}{nm} の光子1個のエネルギーを、eV と J で求めましょう。
\bkjans{\qty{405}{nm}：約\qty{3.06}{eV}（$\qty{4.91e-19}{J}$）。\qty{940}{nm}：約\qty{1.32}{eV}（$\qty{2.12e-19}{J}$）}
\end{jibun}

\label{ph:原-2:ex3}
\begin{phreidai}[LED は1秒に何個の光子を出す？]
波長 \qty{530}{nm} の緑の光を、\qty{1.0}{mW}（1秒に $1.0\times10^{-3}$\,J）出す LED があるとします。1秒に出る光子の数を求めましょう。
\phbamen
ふだんの明かりに、どれほどの数の光子が飛んでいるかを見積もる場面です。
\phmushi
光はすべて \qty{530}{nm} の1つの波長とします（本物の LED は \qty{30}{nm} ほどのはばがある）。
\phzu
光の仕事率 $P=\qty{1.0e-3}{W}$、光子1個 $E=hc/\lambda$、1秒の光子の数 $N$。
\phriyu
1秒に出るエネルギー＝光子の数×光子1個のエネルギー。
\phtateru
\[ P=N\,\frac{hc}{\lambda}\qquad\text{より}\qquad N=\frac{P\lambda}{hc} \]
\phtoku
\[ \frac{hc}{\lambda}=\frac{6.63\times10^{-34}\times3.00\times10^8}{530\times10^{-9}}\approx\qty{3.75e-19}{J} \]
\[ N=\frac{1.0\times10^{-3}}{3.75\times10^{-19}}\approx2.7\times10^{15} \]
\phtashikame
単位：(J/s)÷J＝1/s。\\
大きさ：1秒に約3000兆個。光子1個は小さすぎて、ふだんは光がとぎれなく流れているように見える。
\kotae 1秒に約 $2.7\times10^{15}$ 個
\end{phreidai}
\begin{jibun}
\qty{60}{W}の白熱電球が、電力の 10\,\% を可視光として出すとします（例題1）。可視光を \qty{550}{nm} の光子とみなして、1秒に出る光子の数を求めましょう。
\bkjans{\qty{6.0}{W}÷$\qty{3.62e-19}{J}$ で、1秒に約 $1.7\times10^{19}$ 個}
\end{jibun}

\label{ph:原-2:ex4}
\begin{phreidai}[カリウムに紫の光を当てる]
カリウムの仕事関数を \qty{2.3}{eV} とします。（1）限界の波長を求めましょう。（2）波長 \qty{400}{nm} の光を当てたとき、飛び出す電子の運動エネルギーの最大値、阻止電圧、いちばん速い電子の速さを求めましょう。電子の質量は $\qty{9.11e-31}{kg}$ です。
\phbamen
光電効果で、電子が飛び出すかどうかと、その勢いを調べる場面です。
\phmushi
表面の汚れ・光電面と電極の金属の違いを無視します。相対論は使いません（速さは光の速さの0.2\,\%）。
\phzu
図\ref{ph:原-2:fig-energy}。$W=\qty{2.3}{eV}$、$\lambda=\qty{400}{nm}$。
\phriyu
限界の波長は $hc/\lambda_0=W$。運動エネルギーはアインシュタインの式 $K_0=hf-W$（⑤）。阻止電圧は $eV_0=K_0$。速さは $K_0=\tfrac12mv^2$。
\phtateru
\[ \lambda_0=\frac{hc}{W},\qquad K_0=\frac{hc}{\lambda}-W,\qquad v=\sqrt{\frac{2K_0}{m}} \]
\phtoku
（1）$\lambda_0=1240\div2.3\approx\qty{539}{nm}$。\par
（2）光子は $1240\div400=\qty{3.10}{eV}$ なので、$K_0=3.10-2.3=\qty{0.80}{eV}$。阻止電圧は $V_0=\qty{0.80}{V}$。
\[ v=\sqrt{\frac{2\times0.80\times1.60\times10^{-19}}{9.11\times10^{-31}}}\approx\qty{5.3e5}{m/s} \]
\phtashikame
\qty{400}{nm} は限界の \qty{539}{nm} より短い（光子が大きい）ので、たしかに飛び出す。\\
大きさ：\qty{0.80}{eV} は、\qty{0.80}{V} の電圧で止められる勢い。乾電池1本（\qty{1.5}{V}）より小さい。
\kotae （1）約\qty{540}{nm}。（2）$K_0=\qty{0.80}{eV}$、$V_0=\qty{0.80}{V}$、速さ 約 $\qty{5.3e5}{m/s}$
\end{phreidai}
\begin{jibun}
セシウム（仕事関数 \qty{2.1}{eV} とする）に \qty{450}{nm} の光を当てます。限界の波長、$K_0$、いちばん速い電子の速さを求めましょう。
\bkjans{限界の波長 約\qty{590}{nm}。$K_0=2.76-2.1\approx\qty{0.66}{eV}$。速さ 約 $\qty{4.8e5}{m/s}$}
\end{jibun}

\label{ph:原-2:ex5}
\begin{phreidai}[阻止電圧のグラフから $h$ と仕事関数を]
ある金属に、いろいろな振動数の光を当てて阻止電圧をはかると、表のようになりました（例のための値とします）。グラフにかいて、プランク定数、仕事関数、限界振動数を求めましょう。
\begin{center}\small
\begin{tabular}{c|cccc}
振動数 $f$（$10^{14}$\,Hz） & 5.5 & 6.5 & 7.5 & 8.5 \\ \hline
阻止電圧 $V_0$（V） & 0.27 & 0.69 & 1.10 & 1.51
\end{tabular}
\end{center}
\phbamen
ミリカンと同じ方法で、$h$ をはかる場面です。高校の入試でもよく出る形です。
\phmushi
はかった値の小さな誤差は、直線を引くことでならします。
\phzu
図\ref{ph:原-2:fig-millikan}の実線と点。
\phriyu
$V_0=(h/e)f-W/e$ は直線（⑤）。傾きに $e$ をかけると $h$。直線を延ばして $V_0=0$ になる所が $f_0$、$f=0$ での値が $-W/e$。
\phtateru
両はしの2点を通る直線を引く。
\[ \text{傾き}=\frac{1.51-0.27}{(8.5-5.5)\times10^{14}},\qquad h=e\times(\text{傾き}),\qquad W=hf-eV_0 \]
\phtoku
\[ \text{傾き}=\frac{1.24}{3.0\times10^{14}}\approx\qty{4.13e-15}{V/Hz} \]
\[ h=1.60\times10^{-19}\times4.13\times10^{-15}\approx\qty{6.6e-34}{J.s} \]
仕事関数（eV で書いた数）は、$f=8.5\times10^{14}$\,Hz の点の値から求める。
\[ W=4.13\times10^{-15}\times8.5\times10^{14}-1.51\approx\qty{2.0}{eV} \]
\[ f_0=\frac{2.0}{4.13\times10^{-15}}\approx\qty{4.8e14}{Hz} \]
\phtashikame
$h$ は今の値 $\qty{6.63e-34}{J.s}$ と合う。\\
限界の波長は $c/f_0\approx\qty{620}{nm}$ で、赤い光より短い光で電子が出る金属とわかる。
\kotae $h\approx\qty{6.6e-34}{J.s}$、仕事関数 約\qty{2.0}{eV}、限界振動数 約 $\qty{4.8e14}{Hz}$
\end{phreidai}
\begin{jibun}
別の金属で、振動数 6.0・7.0・8.0・$9.0\times10^{14}$\,Hz のとき、阻止電圧は 0.18・0.60・1.01・\qty{1.42}{V} でした。プランク定数・仕事関数・限界振動数を求めましょう。
\bkjans{両はしの2点の傾きは約 $\qty{4.13e-15}{V/Hz}$。$h\approx\qty{6.6e-34}{J.s}$、仕事関数 約\qty{2.3}{eV}、限界振動数 約 $\qty{5.6e14}{Hz}$}
\end{jibun}

\label{ph:原-2:ex6}
\begin{phreidai}[ケイ素の光センサーが感じる光]
ケイ素の光センサーでは、光子が、電子を禁制帯の幅 \qty{1.12}{eV}（\qty{300}{K}）の分だけもち上げると、電流が流れます（原-8・原-9）。このセンサーが感じる光の、いちばん長い波長を求めましょう。
\phbamen
スマホのカメラや太陽電池が、どんな光まで使えるかを調べる場面です。
\phmushi
光子1個が電子1個をもち上げるとします（熱でもち上がる電子や、2個の光子をいっしょに吸うことを無視）。
\phzu
禁制帯の幅 $E_{\mathrm g}=\qty{1.12}{eV}$、限界の波長 $\lambda_0$。
\phriyu
光電効果の仕事関数と同じように、光子のエネルギーが $E_{\mathrm g}$ 以上なら電子が動ける。限界は $hc/\lambda_0=E_{\mathrm g}$。
\phtateru
\[ \lambda_0=\frac{hc}{E_{\mathrm g}}=\frac{1240\ \mathrm{eV\cdot nm}}{E_{\mathrm g}} \]
\phtoku
\[ \lambda_0=\frac{1240}{1.12}\approx\qty{1110}{nm}=\qty{1.11}{\micro m} \]
\phtashikame
可視光（\qtyrange{0.38}{0.78}{\micro m}）は全部感じ、少し長い赤外線（約 \qty{1.1}{\micro m} まで）も感じる。スマホのカメラが、人には見えないリモコンの赤外線（\qty{940}{nm}）を写せるのはこのため（カメラには、赤外線を弱めるフィルターもついている）。
\kotae 約 \qty{1.11}{\micro m}
\end{phreidai}
\begin{jibun}
ゲルマニウム（禁制帯の幅 \qty{0.66}{eV}）の光センサーが感じる、いちばん長い波長を求めましょう。光ファイバーの通信で使う \qty{1.55}{\micro m} の光を感じますか。
\bkjans{$1240\div0.66\approx\qty{1880}{nm}$（約 \qty{1.9}{\micro m}）。\qty{1.55}{\micro m} は感じる（ケイ素では感じない）}
\end{jibun}

\label{ph:原-2:ex7}
\begin{phreidai}[太陽の光の光子の数と放射圧]
地球の大気の外で、太陽の光に垂直な \qty{1}{m^2} が受ける光は \qty{1361}{W} です（電-15 例題3 と同じ）。光をすべて \qty{500}{nm} の光子とみなして、1秒に \qty{1}{m^2} に当たる光子の数と、全部吸う面が受ける放射圧を求めましょう。
\phbamen
電-15 で電磁波として求めた放射圧を、光子の数と運動量から求めなおす場面です。
\phmushi
太陽の光の波長の広がりを無視し、まん中に近い \qty{500}{nm} で代表します。
\phzu
強さ $I=\qty{1361}{W/m^2}$、光子1個のエネルギー $hc/\lambda$、運動量 $h/\lambda$。
\phriyu
光子の数は（強さ）÷（光子1個のエネルギー）。放射圧は、1秒に当たる光子の運動量の合計（⑥）。
\phtateru
\[ N=\frac{I\lambda}{hc},\qquad p_{\text{圧}}=N\times\frac{h}{\lambda}=\frac{I}{c} \]
\phtoku
\[ \frac{hc}{\lambda}=\frac{6.63\times10^{-34}\times3.00\times10^8}{500\times10^{-9}}\approx\qty{3.98e-19}{J} \]
\[ N=\frac{1361}{3.98\times10^{-19}}\approx3.4\times10^{21}\ \text{個/(m$^2$·s)} \]
\[ p_{\text{圧}}=\frac{1361}{3.00\times10^8}\approx\qty{4.5e-6}{Pa} \]
\phtashikame
放射圧は、波長を何にしても $I/c$ で変わらない（$\lambda$ が消えた）。電-15 の放射圧と同じ答え。\\
大きさ：大気の圧力の200億分の1ほど。全部はね返す帆なら2倍の $\qty{9.1e-6}{Pa}$。
\kotae 1秒に \qty{1}{m^2} あたり約 $3.4\times10^{21}$ 個、放射圧 約 $\qty{4.5e-6}{Pa}$
\end{phreidai}
\begin{jibun}
火星は太陽から地球の1.52倍離れているので、光の強さは $1361\div1.52^2$ になります。火星の近くで、1秒に \qty{1}{m^2} に当たる光子の数（\qty{500}{nm} として）と、全部吸う面の放射圧を求めましょう。
\bkjans{強さ 約 \qty{589}{W/m^2}。光子 約 $1.5\times10^{21}$ 個、放射圧 約 $\qty{2.0e-6}{Pa}$}
\end{jibun}

\label{ph:原-2:ex8}
\begin{phreidai}[光子1個の運動量]
波長 \qty{500}{nm} の光子1個の運動量を求めましょう。電子が同じ運動量をもつときの速さもくらべましょう。
\phbamen
質量のない光子の運動量が、どれくらいかを見る場面です。
\phmushi
とくになし（光子の運動量の式は厳密）。電子は光の速さよりずっと遅いとします。
\phzu
$\lambda=\qty{5.00e-7}{m}$、$p=h/\lambda$、電子 $p=m_{\mathrm e}v$。
\phriyu
光子の運動量は $p=h/\lambda$（⑥）。電子は $p=mv$（力-3）。
\phtateru
\[ p=\frac{h}{\lambda},\qquad v=\frac{p}{m_{\mathrm e}} \]
\phtoku
\[ p=\frac{6.63\times10^{-34}}{5.00\times10^{-7}}\approx\qty{1.3e-27}{kg.m/s} \]
\[ v=\frac{1.33\times10^{-27}}{9.11\times10^{-31}}\approx\qty{1.5e3}{m/s} \]
\phtashikame
$E/c$ で計算しても、例題7の $\qty{3.98e-19}{J}$ を $c$ でわって $\qty{1.3e-27}{kg.m/s}$ で同じ。\\
大きさ：電子なら1秒に \qty{1.5}{km} 進む速さに当たる。光子1個でも、電子を動かせる大きさの運動量をもつ。
\kotae 約 $\qty{1.3e-27}{kg.m/s}$（電子なら約 $\qty{1.5e3}{m/s}$）
\end{phreidai}
\begin{jibun}
波長 \qty{0.10}{nm} の X 線の光子1個の運動量とエネルギー（keV）を求めましょう。
\bkjans{運動量 約 $\qty{6.6e-24}{kg.m/s}$（\qty{500}{nm} の5000倍）、エネルギー $1240\div0.10=\qty{12400}{eV}$、約\qty{12}{keV}}
\end{jibun}

\label{ph:原-2:ex9}
\begin{phreidai}[X 線のコンプトン散乱]
モリブデンの X 線（Kα1 線、波長 \qty{70.9}{pm}）を黒鉛に当てます。$90^\circ$ と $180^\circ$ に散乱した X 線の波長を求めましょう。$h/(m_{\mathrm e}c)=\qty{2.43}{pm}$ です。
\phbamen
コンプトンの実験と同じ X 線で、散乱した X 線の波長のずれを調べる場面です。
\phmushi
黒鉛の中の電子は止まっていて、原子から自由とみます（電子を原子にしばる力は、X 線の光子の \qty{17}{keV} よりずっと小さい）。
\phzu
図\ref{ph:原-2:fig-compton}。$\lambda=\qty{70.9}{pm}$、$\theta=90^\circ$・$180^\circ$。
\phriyu
コンプトン効果の式（⑦）。波長ののびは角だけで決まる。
\phtateru
\[ \lambda'=\lambda+\frac{h}{m_{\mathrm e}c}(1-\cos\theta) \]
\phtoku
$90^\circ$：$\cos90^\circ=0$ なので $\lambda'=70.9+2.43=\qty{73.3}{pm}$。\par
$180^\circ$：$\cos180^\circ=-1$ なので $\lambda'=70.9+2\times2.43=\qty{75.8}{pm}$。
\phtashikame
波長がのびた（エネルギーが減った）ので、光子は電子にエネルギーを渡した。のびは $90^\circ$ で約3\,\%、$180^\circ$ で約7\,\%。
\kotae $90^\circ$：約\qty{73.3}{pm}、$180^\circ$：約\qty{75.8}{pm}
\end{phreidai}
\begin{jibun}
例題9の X 線が $60^\circ$ と $120^\circ$ に散乱したときの波長を求めましょう。
\bkjans{$60^\circ$：$70.9+2.43\times0.5\approx\qty{72.1}{pm}$。$120^\circ$：$70.9+2.43\times1.5\approx\qty{74.5}{pm}$}
\end{jibun}

\label{ph:原-2:ex10}
\begin{phreidai}[γ 線のコンプトン散乱とエネルギー]
セシウム137 が出す \qty{0.662}{MeV} の γ 線が、電子にぶつかって $90^\circ$ と $180^\circ$ に散乱しました。散乱した γ 線と、はじき飛ばされた電子のエネルギーを求めましょう。$m_{\mathrm e}c^2=\qty{0.511}{MeV}$ です。
\phbamen
放射線をはかる検出器の中で起こるぶつかりを、エネルギーで考える場面です（原-13・原-14 で使う）。
\phmushi
電子は止まっていて自由とみます。
\phzu
$E=\qty{0.662}{MeV}$、散乱後 $E'$、電子 $E-E'$。
\phriyu
⑦のエネルギーの形の式。電子のエネルギーは、エネルギーの保存から差。
\phtateru
\[ E'=\frac{E}{1+\dfrac{E}{m_{\mathrm e}c^2}(1-\cos\theta)} \]
\phtoku
$E/(m_{\mathrm e}c^2)=0.662\div0.511\approx1.296$。\par
\[ 90^\circ:\ E'=\frac{0.662}{1+1.296}\approx\qty{0.288}{MeV},\qquad\text{電子}\ \ 0.662-0.288=\qty{0.374}{MeV} \]
\[ 180^\circ:\ E'=\frac{0.662}{1+2\times1.296}\approx\qty{0.184}{MeV},\qquad\text{電子}\ \ \qty{0.478}{MeV} \]
\phtashikame
$180^\circ$ のときが、電子がもらえるいちばん大きいエネルギーで、これより大きいエネルギーの電子はできない。検出器でエネルギーをはかると、\qty{0.478}{MeV} の所で数がとぎれる（コンプトン端：原-13）。
\kotae $90^\circ$：γ 線 \qty{0.288}{MeV}、電子 \qty{0.374}{MeV}。$180^\circ$：γ 線 \qty{0.184}{MeV}、電子 \qty{0.478}{MeV}
\end{phreidai}
\begin{jibun}
PET の検査で使う \qty{0.511}{MeV} の γ 線が、$90^\circ$ と $180^\circ$ に散乱したときの、散乱した γ 線と電子のエネルギーを求めましょう。
\bkjans{$E/(m_{\mathrm e}c^2)=1$。$90^\circ$：γ 線・電子とも約\qty{0.256}{MeV}。$180^\circ$：γ 線 約\qty{0.170}{MeV}、電子 約\qty{0.341}{MeV}}
\end{jibun}

\label{ph:原-2:ex11}
\begin{phreidai}[モード1つは、どれだけのエネルギーをもつ？]
2けたで書きます。波長 \qty{550}{nm} のモード1つの平均のエネルギーは、\qty{300}{K}（室温）と \qty{5772}{K}（太陽の表面）で、$k_{\mathrm B}T$ の何倍ですか。波だけで考えたとき（いつも $k_{\mathrm B}T$）とくらべましょう。
\phbamen
紫外の破綻がなぜ起こらないかを、数で確かめる場面です。
\phmushi
とくになし（⑩の式は、熱のつり合いにある光について厳密）。
\phzu
$x=hc/(\lambda k_{\mathrm B}T)$、$\langle\varepsilon\rangle/(k_{\mathrm B}T)=x/(e^x-1)$。
\phriyu
モード1つの平均のエネルギーの式（⑩）を $k_{\mathrm B}T$ でわると $x/(e^x-1)$。
\phtateru
\[ x=\frac{hc}{\lambda k_{\mathrm B}T},\qquad\frac{\langle\varepsilon\rangle}{k_{\mathrm B}T}=\frac{x}{e^x-1} \]
\phtoku
光子1個は $hc/\lambda=\qty{2.25}{eV}$。
\[ \qty{300}{K}:\ x=\frac{2.25}{0.0259}\approx87,\qquad\frac{x}{e^x-1}\approx87\times e^{-87}\approx1.2\times10^{-36} \]
\[ \qty{5772}{K}:\ k_{\mathrm B}T=\qty{0.497}{eV},\quad x\approx4.53,\qquad\frac{x}{e^x-1}=\frac{4.53}{e^{4.53}-1}\approx0.049 \]
\phtashikame
\qty{300}{K}の部屋の壁は、緑の光をまったくと言っていいほど出さない（だから暗い部屋は暗い）。波だけで考えると、どちらも1（$k_{\mathrm B}T$）になってしまい、部屋の壁が光りかがやくことになる。
\kotae \qty{300}{K}：約 $1.2\times10^{-36}$ 倍、\qty{5772}{K}：約0.049倍（波だけの考えでは、どちらも1倍）
\end{phreidai}
\begin{jibun}
\qty{300}{K}で、波長 \qty{10}{\micro m}（赤外線）と \qty{1}{mm}（電波）のモード1つの平均のエネルギーは、$k_{\mathrm B}T$ の何倍ですか。
\bkjans{\qty{10}{\micro m}：$x\approx4.8$、約0.040倍。\qty{1}{mm}：$x\approx0.048$、約0.98倍（波だけの考えとほぼ同じ）}
\end{jibun}

\label{ph:原-2:ex12}
\begin{phreidai}[熱-13 の定数を、プランク定数から]
4けたで書きます。3つの定数を
\[ h=6.62607\times10^{-34}\ \mathrm{J\cdot s},\qquad c=2.99792\times10^{8}\ \mathrm{m/s} \]
\[ k_{\mathrm B}=1.38065\times10^{-23}\ \mathrm{J/K} \]
として、⑪の式から $c_1$・$c_2$・$\sigma$・ウィーンの定数 $b_{\mathrm W}$ を求め、熱-13 の値とくらべましょう。
\phbamen
熱-13 で「測定で決まる」とした定数が、3つの基本の定数で書けることを確かめる場面です。
\phmushi
とくになし。
\phzu
熱-13 の値は、$c_1=3.742\times10^8$\,W·\unit{\micro m}$^4$/m$^2$、$c_2=\num{14388}$\,\unit{\micro m}·K、$\sigma=5.670\times10^{-8}$\,W/(m$^2$·K$^4$)、ウィーンの定数 \num{2898}\,\unit{\micro m}·K。
\phriyu
⑪で導いた式。
\phtateru
\[ c_1=2\pi hc^2,\quad c_2=\frac{hc}{k_{\mathrm B}},\quad\sigma=\frac{2\pi^5k_{\mathrm B}^4}{15h^3c^2},\quad b_{\mathrm W}=\frac{hc}{4.965\,k_{\mathrm B}} \]
\phtoku
\[ c_1=2\pi\times6.62607\times10^{-34}\times(2.99792\times10^8)^2\approx\qty{3.742e-16}{W.m^2} \]
\[ c_2=\frac{6.62607\times10^{-34}\times2.99792\times10^8}{1.38065\times10^{-23}}\approx\qty{1.439e-2}{m.K} \]
同じように、$\sigma\approx\qty{5.670e-8}{W/(m^2.K^4)}$、$b_{\mathrm W}=c_2\div4.965\approx\qty{2.898e-3}{m.K}$。
\phtashikame
\qty{1}{m^4} は $10^{24}$\,\unit{\micro m}$^4$ なので、$c_1$ は $3.742\times10^8$\,W·\unit{\micro m}$^4$/m$^2$。$c_2$ は $1.439\times10^{4}$\,\unit{\micro m}·K。どちらも熱-13 の値と同じ。
\kotae $c_1$ 約 $\qty{3.742e-16}{W.m^2}$、$c_2$ 約 $\qty{1.439e-2}{m.K}$、$\sigma$ 約 $\qty{5.670e-8}{W/(m^2.K^4)}$、$b_{\mathrm W}$ 約 $\qty{2.898e-3}{m.K}$（熱-13 と同じ）
\end{phreidai}
\begin{jibun}
3けたに丸めた $h=\qty{6.63e-34}{J.s}$、$c=\qty{3.00e8}{m/s}$、$k_{\mathrm B}=\qty{1.38e-23}{J/K}$ を使って $\sigma$ を計算すると、いくらになりますか。正しい値と何\,\% 違いますか。違いが大きくなるわけも答えましょう。
\bkjans{約 $\qty{5.64e-8}{W/(m^2.K^4)}$、約0.5\,\% 小さい。$\sigma$ は $k_{\mathrm B}^4/h^3$ に比例するので、$k_{\mathrm B}$ の丸めのずれが4倍、$h$ のずれが3倍になって効くから}
\end{jibun}

%% ======================== まちがえやすいところ ========================
\Needspace{13\baselineskip}\bksection{tsumazuki}{まちがえやすいところをチェック}
\begin{tsumazuki}{強い光なら、どんな色でも電子が飛び出す？}
\bkng{赤い光でも、うんと強くすれば、亜鉛から電子が飛び出す。}
\bkok{光子1個のエネルギー $hf$ が仕事関数より小さいと、光の強さ（光子の数）を増やしても飛び出さない。}
光子は、ふつう1個ずつ電子に吸われます（④）。とても強いレーザーだけが例外です（⑤のモデル化の札）。
\end{tsumazuki}

\begin{tsumazuki}{振動数を2倍にすると、電子の運動エネルギーも2倍？}
\bkng{光の振動数を2倍にすると、$K_0$ も2倍。}
\bkok{$K_0=hf-W$ なので、$f$ を2倍にすると $K_0$ は2倍より大きく増える。比例するのは $K_0+W$。}
例題4の \qty{400}{nm}（$K_0=\qty{0.80}{eV}$）を \qty{200}{nm}（$hf=\qty{6.2}{eV}$）にすると、$K_0=\qty{3.9}{eV}$ で約5倍です。
\end{tsumazuki}

\begin{tsumazuki}{コンプトン効果の波長ののびは、もとの波長で変わる？}
\bkng{波長の長い光ほど、散乱で波長が大きくのびる。}
\bkok{のび $\Delta\lambda=\lambda_{\mathrm C}(1-\cos\theta)$ は、もとの波長によらない。}
同じのびでも、もとの波長が短いほど、のびの割合（とエネルギーの減り方）が大きくなります（⑦）。
\end{tsumazuki}

\begin{tsumazuki}{光の山の波長と、山の振動数は同じ光？}
\bkng{\qty{5772}{K}の光の山は \qty{0.50}{\micro m} だから、振動数で見た山は、それを振動数に直した $\qty{6.0e14}{Hz}$。}
\bkok{振動数で分布をかくと、山は $hf=2.82k_{\mathrm B}T$ の所で、$\qty{3.4e14}{Hz}$（波長に直すと \qty{0.88}{\micro m}）。}
「波長のはば \qty{1}{nm} ごと」と「振動数のはば \qty{1}{Hz} ごと」では、同じはばに入る光が違うので、山の位置が変わります（⑪で $\dd f$ を $\dd\lambda$ に書きかえたとき、$c/\lambda^2$ がかかった）。
\end{tsumazuki}

%% ======================== 工学ではこう使う ========================
\Needspace{10\baselineskip}\bksection{phgear}{工学ではこう使う}
\begin{phkougaku}[光を数える・温度をはかる]
\textbf{光をはかる部品を選ぶ}：光センサーを選ぶときは、はかりたい光の光子1個のエネルギーと、センサーの限界（光電面の仕事関数、半導体の禁制帯の幅）をくらべます。紫外線から可視光には光電子増倍管やケイ素のセンサー、\qty{1.1}{\micro m} より長い赤外線にはゲルマニウムなどの禁制帯のせまい半導体を使います（例題6）。とどいた光子のうち、電子を出すものの割合を\textbf{量子効率}といい、センサーの良さを表す数です。

\textbf{さわらずに温度をはかる}：放射温度計は、物が出す赤外線の強さを、プランクの式（⑪）とくらべて温度を求めます。物は黒体ではないので、放射率（熱-13）を正しく入れることが大切です。溶けた鉄（約 \qty{1800}{K}）のような高い温度は、2つの波長の強さの比から求める方法（2色温度計）も使われます。

\begin{phketa}[光と光子の数の目安]
\phketaline{$hc$}{\qty{1240}{eV.nm}（eV＝1240÷nm）}
\phketaline{可視光の光子}{\qtyrange{1.6}{3.3}{eV}}
\phketaline{金属の仕事関数}{\qtyrange{2}{5}{eV}}
\phketaline{\qty{1}{mW}の緑の光の光子}{1秒に約 $3\times10^{15}$ 個}
\phketaline{電子のコンプトン波長}{\qty{2.43}{pm}}
\phboxhead{答えが桁で合っているか}
可視光の光子が \qty{1}{eV} より小さい、または \qty{10}{eV} より大きく出たら、nm と m の換算、J と eV の換算を見直しましょう。
\end{phketa}
\end{phkougaku}
\phstation{光エレクトロニクス,計測工学,放射測温,量子光学}{光電面、量子効率、仕事関数、阻止電圧、光子計数、光電子増倍管、黒体炉、放射率、コンプトン散乱}

%% ======================== 章末問題 ========================
\bkendproblems
\bkproblemgroup{基本}{3}
\bkproblem 波長 \qty{620}{nm} と \qty{310}{nm} の光子1個のエネルギーを、eV と J で求めましょう。
\bkpans{\qty{620}{nm}：\qty{2.0}{eV}（$\qty{3.2e-19}{J}$）。\qty{310}{nm}：\qty{4.0}{eV}（$\qty{6.4e-19}{J}$）}
\bkproblem 仕事関数 \qty{2.4}{eV} の金属に、\qty{300}{nm} の紫外線を当てました。（1）限界の波長を求めましょう。（2）飛び出す電子の運動エネルギーの最大値と、阻止電圧を求めましょう。
\bkpans{（1）$1240\div2.4\approx\qty{520}{nm}$。（2）$4.13-2.4\approx\qty{1.7}{eV}$、阻止電圧 約\qty{1.7}{V}}
\bkproblem 波長 \qty{620}{nm} の赤い光を \qty{3.0}{mW} 出す LED があるとします。1秒に出る光子の数を求めましょう。
\bkpans{光子1個 $\qty{3.2e-19}{J}$。$3.0\times10^{-3}\div(3.2\times10^{-19})$ で、約 $9.4\times10^{15}$ 個}
\bkproblem \qty{0.662}{MeV} の γ 線が電子にぶつかって $120^\circ$ に散乱しました。散乱した γ 線と、電子のエネルギーを求めましょう。
\bkpans{$1-\cos120^\circ=1.5$ を入れて、散乱した γ 線は約\qty{0.225}{MeV}、電子は約\qty{0.437}{MeV}}
\bkproblem （1）\qty{3000}{K}の黒体の光の山の波長を求めましょう。（2）温度を4倍にすると、山の波長と、出る光全体の強さは何倍になりますか。
\bkpans{（1）$2.898\times10^{-3}\div3000\approx\qty{0.97}{\micro m}$。（2）山の波長は $\tfrac14$ 倍、全体の強さは $4^4=256$ 倍}
\bkproblem 光電管で阻止電圧が \qty{1.5}{V} でした。飛び出した電子のうち、いちばん速いものの速さを求めましょう。
\bkpans{$K_0=\qty{1.5}{eV}=\qty{2.4e-19}{J}$。$v=\sqrt{2K_0/m_{\mathrm e}}\approx\qty{7.3e5}{m/s}$}
\bkproblem 宇宙マイクロ波背景放射は、温度 \qty{2.725}{K} の黒体の光です。山の波長と、その波長の光子1個のエネルギー（eV）を求めましょう。
\bkpans{山の波長 $2.898\times10^{-3}\div2.725\approx\qty{1.06e-3}{m}$（約\qty{1.06}{mm}）。光子1個 $1240\div(1.06\times10^{6})\approx\qty{1.2e-3}{eV}$}
\bkproblemgroup{チャレンジ}{5}
\bkproblem ある金属に \qty{300}{nm} の光を当てると阻止電圧は \qty{2.03}{V}、\qty{400}{nm} では \qty{1.00}{V} でした。この2つのデータだけから、プランク定数と仕事関数を求めましょう。
\bkpans{振動数は $\qty{1.00e15}{Hz}$ と $\qty{7.5e14}{Hz}$。2点を結ぶ直線の傾きは約 $\qty{4.12e-15}{V/Hz}$ で、$h\approx\qty{6.6e-34}{J.s}$。仕事関数は約\qty{2.1}{eV}}
\bkproblem \qty{300}{K}の黒体で、波長 \qty{30}{\micro m} のモード1つの平均のエネルギーは、波だけで考えたとき（$k_{\mathrm B}T$）の何倍ですか。また、波長がこれより長い光と短い光で、レイリー・ジーンズの式はどちらでよく合うか答えましょう。
\bkpans{$x=hc/(\lambda k_{\mathrm B}T)\approx1.60$、$x/(e^x-1)\approx0.41$ 倍。$x$ が小さい（波長が長い）ほど1に近づくので、長い波長でよく合う}
\bkproblem 白熱電球のフィラメントの温度を \qty{2800}{K} から \qty{3000}{K} に上げると、出る光全体の強さは何倍になりますか。例題1のプログラムで \qty{3000}{K}の可視光の割合を求め、出る可視光の量が何倍になるかを求めましょう（フィラメントは黒体とし、表面の大きさは同じとする）。
\bkpans{全体は $(3000/2800)^4\approx1.32$ 倍。可視光の割合は約 9.8\,\% から約 12.7\,\% に。可視光の量は $1.32\times12.7\div9.8\approx1.7$ 倍（温度を少し上げると、明かりの効率がよくなる。ただし、フィラメントは蒸発しやすくなり寿命が短くなる）}

\begin{furikaeri}
\begin{checklist}
\item 熱した物の色が温度で変わるわけを、光の分布の山の動きで説明できる。\bkfurref{ph:原-2:sec1}
\item 光子のエネルギー $E=hf=hc/\lambda$ を、eV と J で求められる。\bkfurref{ph:原-2:sec3}
\item 光電効果のきまりを光子で説明し、$K_0=hf-W$ と阻止電圧を使える。\bkfurref{ph:原-2:sec5}
\item 光子の運動量 $p=h/\lambda$ から、放射圧を求められる。\bkfurref{ph:原-2:sec6}
\item コンプトン効果の式で、散乱した X 線・γ 線の波長とエネルギーを求められる。\bkfurref{ph:原-2:sec7}
\item モードの数と平均のエネルギーから、プランクの式・ウィーンの変位則・$\sigma$ を導ける。\bkfurref{ph:原-2:sec11}
\end{checklist}
\end{furikaeri}
\bkfukidashi{mina}{光が波でもあり、つぶでもあるなんて、ふしぎ！ 電子も、つぶなのに波のようにふるまうのかな？}
